\documentclass[12pt]{article}


\usepackage{amsmath,amssymb,amsfonts}
\usepackage{amsthm}
\usepackage{bm}
\usepackage{geometry}
\usepackage[numbers,super,sort&compress]{natbib}
\usepackage{booktabs}
\usepackage{multirow}
\geometry{margin=1in}


% Hyperlinks (invisible style: normal text color; no boxes)
\usepackage[hidelinks]{hyperref}
\hypersetup{linktoc=all}
\usepackage[capitalize,noabbrev,nameinlink]{cleveref}


\newtheorem{definition}{Definition}[section]
\newtheorem{assumption}[definition]{Assumption}
\newtheorem{remark}[definition]{Remark}
\newtheorem{proposition}[definition]{Proposition}
\newtheorem{example}[definition]{Example}




\begin{document}


% ==================================================================
% STYLE GUIDELINES
% ==================================================================
% 1. State object and claim once; no meta-commentary or repetition.
% 2. Lead with operational target (what to compute/test), then formalism.
% 3. Humility is informational: scope conditions and failure modes, not hedging.
% 4. Compact orientation aids (one toy example; short roadmap), then move on.
% 5. Empirical tests specify: what we compute, expected result, implications.
% 6. Timelessness: no version numbers in prose; frame progress as cumulative.
% 7. Nulls as constraints: non-findings are model-space restrictions, not failures.
% ==================================================================


\title{Embedding-Informed Task Distance for Occupational Mobility}

\maketitle


\begin{abstract}
We develop a modular task-space framework mapping $(\mathcal{T}, \mathcal{I}, \mathcal{S}, \mathcal{M}) \to (\Delta\boldsymbol{\rho}, \Delta \mathbf{L}, \Delta \mathbf{W})$: from task representation, institutional structure, shock profiles, and adjustment mechanisms to occupation-specific changes in task distributions, employment, and wages.

\textbf{Module status.} $\mathcal{T}$: Validated---embedding-informed distance achieves 13.8--14.1\% pseudo-$R^2$ on 89,329 CPS transitions; the embedding ground metric (+83\% vs.\ identity) drives improvement, not distributional treatment (Wasserstein adds nothing beyond centroid). We find that most of the predictive gain comes from the embedding representation of inter-task similarity, not from the distributional treatment of occupations. Embedding-based specifications achieve pseudo-$R^2$ of 13.8--14.1\%, compared to 6.1--8.1\% for non-embedding structured-vector baselines. $\mathcal{I}$: Validated---institutional distance adds predictive power conditional on $\mathcal{T}$ ($t = 33.7$). $\mathcal{S}$: Validated---external AI exposure scores are orthogonal to geometry ($r = -0.02$); geometry substantially outperforms historical baselines ($\Delta\text{LL} = +23{,}879$ out-of-period). $\mathcal{M}$: Calibrated externally (8.74 wage-years per unit centroid distance).

\textbf{Key findings.} Aggregate flows are demand-dominated ($\rho = 0.80$); geometry captures supply-side feasibility ($\rho \approx 0.12$ per-origin). Task-distance constraints are structurally stable across pre/post COVID ($\Delta\alpha < 1\%$, $p = 0.72$), with selectively elevated hiring standards for teleworkable occupations ($\delta_4 = -0.086$, $p = 0.01$).

\end{abstract}


\vspace{1em}
\noindent\fbox{\parbox{\dimexpr\textwidth-2\fboxsep-2\fboxrule}{%
\textbf{Authorship \& AI Contribution}\\[0.5em]
This research was conducted using AI-assisted methods throughout. The initial conceptualization (task-space framing, modular architecture) is the author's; theory development, experimental design, code implementation, statistical analysis, and writing were performed by large language models (Claude Opus 4.5, GPT-5.2) based on the author's direction and iterative feedback. All reported results are reproducible via the provided codebase; methodology and interpretation have not yet been reviewed by a domain expert.\\[0.3em]
\textit{Reproducibility:} Code repository available at \texttt{github.com/spencer2718/task-space-model}. Runtime: $\sim$45 minutes on standard hardware.
}}
\vspace{1em}


\begin{center}
\textit{This document specifies an ongoing research program. The architecture is viable; modules can be extended incrementally.}
\end{center}


% ==================================================================
% @AGENT: ORIENTATION SECTION
% Contains paper framing and contribution summary.
% Update contribution list when major findings change.
% Current: Module validation + supply-demand decomposition framing
% ==================================================================

\section{Introduction}
\label{sec:introduction}


Technological change acts on tasks. Within any occupation, technology automates some activities, augments others, and enables new ones. These task-level changes---shifts in what workers actually do---propagate to occupation-level outcomes: employment and wages adjust as the demand for different task bundles evolves. Technology acts on tasks; occupations are distributions over tasks; labor market outcomes are downstream aggregations. Understanding labor market responses to technology therefore requires modeling how task distributions change, where workers with given skill profiles can feasibly transition, and how wages adjust in equilibrium.


This framing, inspired by the task-based approach to labor markets, clarifies what a complete analysis requires. Existing research addresses pieces: routine task intensity measures characterize which tasks are automatable; AI exposure indices identify which occupations contain vulnerable tasks; gravity models characterize worker flows between occupations. But no framework integrates these components---task representation, institutional barriers, shock characterization, and adjustment mechanisms---into a unified architecture for analyzing labor market adjustment.


We propose a modular task-space framework that targets complete structural mapping from task representations, technology shocks, and institutions to equilibrium labor market outcomes. The framework is an organizing architecture for a research program, not a deliverable model. Its value is taxonomic: it clarifies which components a complete analysis requires, enabling systematic validation and incremental progress.

The framework has four modules:
\begin{itemize}
    \item $\mathcal{T}$ \textbf{(Task representation):} How to encode occupation skill content. We validate embedding-informed distance on O*NET task embeddings as the primary measure.
    \item $\mathcal{I}$ \textbf{(Institutional structure):} Non-skill barriers to mobility---credentials, job zones, licensing. We validate that institutional distance adds predictive power conditional on semantic distance.
    \item $\mathcal{S}$ \textbf{(Shock profiles):} Which capabilities are affected by technology. We integrate external measures (AI Occupational Exposure) rather than estimating endogenously.
    \item $\mathcal{M}$ \textbf{(Mechanisms):} How shocks propagate to outcomes---switching costs, search frictions, equilibrium adjustment. We identify data requirements; full estimation is future work.
\end{itemize}

A complete framework would map $(\mathcal{T}, \mathcal{I}, \mathcal{S}, \mathcal{M}) \to (\Delta\boldsymbol{\rho}, \Delta \mathbf{L}, \Delta \mathbf{W})$: from inputs to occupation-specific changes in task distributions, employment, wages. No such framework exists. This paper validates candidate specifications for $\mathcal{T}$ and $\mathcal{I}$, integrates external $\mathcal{S}$, and identifies what $\mathcal{M}$ requires.


The paper's contributions are framed as module validation checkpoints, with our primary contribution being \emph{the 2$\times$2 decomposition and the choice-model framework} rather than conceptual innovation. Prior work has validated occupational similarity against transition data: \citet{frank2024network} show O*NET skill similarity predicts CPS transitions, and \citet{dawson2021skill} validate skill-based occupation distance against Australian survey transitions. Embedding approaches (e.g., occ2vec, JobBERT) typically validate against O*NET's own similarity measures or expert ratings---potentially circular targets. We validate against revealed worker behavior: 89,329 actual CPS transitions, using a factorial design that isolates which design choice---representation or aggregation---drives predictive improvement.

First, the $\mathcal{T}$ module is validated through a 2$\times$2 methodology comparison isolating embedding choice from aggregation method. Embedding-based distances (using sentence embeddings of Detailed Work Activities) achieve 13.8--14.1\% pseudo-$R^2$ for predicting individual transitions, compared to 6--8\% for non-embedding O*NET structured-vector baselines---a 74.9\% relative improvement. The mechanism is the embedding ground metric: semantic task similarity captures that ``operating forklift'' $\approx$ ``driving delivery vehicle'' even as distinct task labels. Wasserstein optimal transport provides theoretical grounding as minimum semantic transformation cost; simpler centroid-averaging methods approximate this geometry closely ($\rho = 0.95$ correlation in distance rankings) and marginally outperform Wasserstein after correcting for diagonal bias. The contribution is the 2$\times$2 factorial design isolating the embedding representation from the aggregation method, and the validation against a large US transition sample in a choice-model framework. The finding that most of the predictive gain comes from the embedding representation, not the distributional treatment, has not been established in prior work.

Second, the $\mathcal{I}$ module is validated: institutional distance (job zones plus certification requirements) adds predictive power conditional on embedding-based semantic distance ($\beta = 0.139$, $t = 33.7$). The separability assumption holds---skill barriers and credential barriers are distinct friction sources. Directional asymmetry (upward vs.\ downward barriers) is sample-dependent: ratio varies from 0.06 to 2.79 across specifications. Neither pure ``credential gate'' nor pure ``symmetric friction'' theories are universally supported; this heterogeneity constrains future $\mathcal{I}$ specifications.

Third, $\mathcal{S}$ integration is validated: external AI Occupational Exposure scores \citep{FeltenRajSeamans2021} are orthogonal to embedding geometry ($r = -0.02$). Shock profiles identify which occupations contain exposed tasks; geometry identifies \emph{where} exposed workers can feasibly transition. The orthogonality confirms these are complementary constructs. Geometry substantially outperforms historical baselines in probabilistic scoring ($\Delta \text{log-likelihood} = +23{,}879$).

Fourth, $\mathcal{M}$ module requirements are identified. Endogenous switching cost estimation fails with our data: occupation-mean wages yield negative coefficients, indicating that aggregate wages do not capture the entry wages switchers actually receive. External calibration to structural estimates \citep{DixCarneiro2014} yields 8.74 wage-years per unit centroid distance. This is an interim approach; full $\mathcal{M}$ estimation requires individual-level wage data at transition.

Fifth, scope is empirically quantified through supply-demand decomposition. Aggregate occupational inflows are demand-dominated: job openings alone predict aggregate inflows with $\rho = 0.80$, while geometry alone achieves $\rho \approx 0.00$. For pathway direction at the origin level, geometry provides modest but statistically significant information (per-origin $\rho \approx 0.12$, $n = 233$ origins). This decomposition validates the framework architecture's separation of supply-side feasibility (where workers \emph{can} go given skill geometry) from demand-side realization (where they \emph{do} go given labor market conditions). The framework measures structural feasibility---the supply-side input to equilibrium analysis.

Sixth, a gravity model analysis confirms that task distance explains a modest share of aggregate bilateral flow variation ($\sim$2.5\% partial $R^2$), consistent with \citet{CortesGallipoli2018}'s finding that task-specific costs account for no more than 15\% of total transition costs across most occupation pairs. The divergence between individual choice prediction (where embeddings excel) and aggregate flow prediction (where simpler measures are competitive) reflects different margins: individual workers navigating among connected occupations use fine-grained task similarity, while aggregate patterns depend primarily on whether occupation pairs exchange workers at all.

Seventh, structural stability is tested via pre/post COVID comparison. Aggregate task-distance penalties are invariant across periods (coefficient change $< 1\%$), supporting the interpretation that embedding-based distance captures structural skill transferability rather than contingent revealed preferences. Heterogeneity analysis reveals that teleworkable occupations experienced selectively elevated hiring standards post-COVID---consistent with \citet{ModestinoShoagBallance2020}'s framework where expanded applicant pools enable heightened employer selectivity.


The paper explores a model space defined by three axes: (i) \emph{representation}---how to encode tasks; (ii) \emph{geometry}---how to measure occupation distance; (iii) \emph{shocks}---how to characterize technology. Each empirical test constrains this space. The 2$\times$2 methodology comparison resolves representation: embedding-informed distance substantially outperforms non-embedding structured-vector baselines. The RTI construct validity test establishes that semantic geometry is orthogonal to routine task intensity ($r = -0.06$)---the geometry captures mobility friction, not automation susceptibility. These are distinct constructs requiring distinct measurement.


Our approach complements structural estimation \citep{Traiberman2019} by trading parametric richness for applicability when administrative data is unavailable. We extend the structured-vector task-distance tradition \citep{GathmannSchonberg2010,CortesGallipoli2018,MouwEtAl2024,SchuStanTaska2024} with embedding-informed geometry, and we extend embedding approaches \citep{ZhangEtAl2021,DecorteEtAl2022} with validation against revealed mobility rather than O*NET self-similarity. The supply-demand decomposition ($\rho \approx 0.80$ demand, $\rho \approx 0.12$ geometry) is consistent with \citet{CortesGallipoli2018}'s finding that task-related costs account for no more than 15\% of total transition costs. The framework provides infrastructure for cumulative progress: validated modules can be reused as new capabilities are added.


\Cref{sec:literature} reviews related literatures organized by framework module. \Cref{sec:theory} develops the framework architecture and module specifications. \Cref{sec:data} operationalizes each component using O*NET data. \Cref{sec:validation} reports module validation results including geometry comparison, CPS mobility, shock integration, cost calibration, supply-demand decomposition, and aggregate flow analysis. \Cref{sec:extensions} documents extensions, completed investigations, and deferred work. \Cref{sec:limitations} documents scope conditions. \Cref{sec:applications} discusses policy applications and the research roadmap. \Cref{sec:conclusion} concludes with module status and next steps.


\section{Literature Review}
\label{sec:literature}


This section reviews literatures that inform framework module development, organized by the component each addresses.


\subsection{Task Representation}
\label{subsec:lit-task}


The routine-biased technological change paradigm established that technology acts on tasks, not occupations directly \citep{AutorLevyMurnane2003}. Computerization disproportionately substitutes for routine, rule-based tasks while complementing non-routine analytic and interactive tasks. Subsequent work formalized task assignment and emphasized that equilibrium implications for employment and wages cannot be recovered from skill aggregates alone \citep{AcemogluAutor2011}.

Our framework is inspired by the task-based approach to technological change \citep{AcemogluAutor2011,AcemogluRestrepo2018}, which emphasizes that technology acts on tasks and labor market outcomes are aggregations of task-level effects. We do not implement their specific model of task creation and destruction; rather, we build measurement infrastructure---task geometry, institutional distance, shock integration---that could support such models.

A broad literature represents occupations using structured descriptor vectors from O*NET, DOT, or related task datasets and computes pairwise proximity mechanically using correlation-, cosine-, or distance-based measures. Implementations differ in descriptors, weighting schemes, and estimation frameworks. The canonical empirical methodology constructs task measures from the Dictionary of Occupational Titles using specific variables aggregated into routine, abstract, and manual indices \citep{AutorDorn2013}. O*NET has become the primary measurement infrastructure, with 2,087 Detailed Work Activities linked to 894 O*NET-SOC occupations with DWA data. More recent work has brought structured-vector distances into worker-level estimation: \citet{MouwEtAl2024} use O*NET skill similarity inside a conditional-logit destination-choice model; \citet{Macaluso2025} uses O*NET skill remoteness to predict post-layoff labor market outcomes; \citet{SchuStanTaska2024} use O*NET task distance to analyze determinants of occupational transition shares; \citet{CarrilloTudelaEtAl2025} validate O*NET distance against corrected CPS transitions; and \citet{KudlyakWolcott2019} construct occupation skill-vector distances from O*NET profiles. Our non-embedding baselines belong to this tradition: importance-weighted O*NET DWA vectors evaluated with cosine and Euclidean distance.

Text embedding approaches map technologies to occupations using semantic similarity. \citet{Webb2020} matches patent text to job descriptions (cross-domain matching). The present paper matches task text to task text within O*NET (intra-domain matching), avoiding challenges of aligning heterogeneous corpora. \citet{EloundouEtAl2024} provide task-level AI exposure indices using GPT-4 capability ratings.

Applying embeddings to occupational descriptions is not new. \citet{ZhangEtAl2021} (occ2vec) embed occupation titles and validate against O*NET similarity measures, achieving 53\% correlation between occupations and 77\% within occupations. \citet{DecorteEtAl2022} (JobBERT) fine-tune BERT on ESCO skill labels for job-resume matching. \citet{HampoleEtAl2025} use GTE-Large embeddings for AI exposure assessment, computing cosine similarity between AI applications and O*NET task descriptions. These approaches typically validate against O*NET's own similarity measures or expert ratings---potentially circular, since O*NET structure influences embedding geometry.

Three additional precedents merit discussion. \citet{dawson2021skill} build occupation distance bottom-up from skill-level similarity using co-occurrence in Australian job ads, then validate against HILDA survey transitions using XGBoost classification. Their feature ablation shows skill-distance is the most important predictor. They do not use neural text embeddings or study aggregation operators; their similarity is co-occurrence-based (Revealed Comparative Advantage), and their evaluation is binary classification (76\% accuracy) against random negative counterfactuals rather than discrete-choice likelihood. \citet{frank2024network} validate O*NET-derived skill similarity against 36,536 CPS transitions (2012--2018), showing skill similarity predicts observed transition rates. They use structured O*NET skill importance ratings (not text embeddings) and do not compare representation choices or aggregation methods. O*NET's own related-occupations pipeline \citep{onet2024related} now uses SBERT embeddings of task statements and DWAs for computing occupational relatedness, but this has not been validated against transition data.

Relative to this literature, the contribution is to compare semantically enriched embedding-based representations against transparent structured-vector baselines in a unified destination-choice setting, using a factorial design that isolates which design choice---representation or aggregation---drives predictive improvement. Prior structured-vector work \citep{MouwEtAl2024,SchuStanTaska2024} validates O*NET distance against transitions but does not compare representation choices or aggregation methods. The finding that most of the predictive gain comes from the embedding representation, not the distributional treatment, has not been established in prior work.

Our contribution is methodological: validating embedding-based distance against \emph{revealed} worker behavior (89,329 CPS transitions) rather than O*NET self-similarity. The 2$\times$2 design isolates what matters: the embedding ground metric, not the aggregation method, drives the 74.9\% improvement over non-embedding structured-vector baselines. This validation target---actual mobility patterns---addresses circularity concerns and quantifies practical predictive gains.

Optimal transport in economics provides theoretical foundation for principled distance metrics \citep{Galichon2016}. The Wasserstein ``earth mover's'' interpretation---minimum cost to transform one skill portfolio into another---is economically intuitive. We find that simpler centroid-based methods approximate Wasserstein geometry closely ($\rho = 0.95$) and marginally outperform after correcting for diagonal bias, suggesting that for mobility prediction, the embedding representation matters more than the distributional treatment.


\subsection{Institutional Barriers}
\label{subsec:lit-institutional}


The licensing literature documents substantial effects on labor mobility. \citet{Jackson2023} finds licensing reduces entry probability by 24\% with no significant effect on exit---consistent with one-way gate structure. \citet{KleinerXu2020} find licensing reduces \emph{both} entry and exit rates by similar magnitudes, suggesting symmetric friction rather than gates. \citet{BlairChung2019} estimate licensing reduces labor supply by 17--27\%. The literature remains unsettled on whether credentials function as gates or friction.

Human capital specificity creates additional barriers. \citet{GathmannSchonberg2010} find up to 52\% of wage growth derives from task-specific human capital destroyed in any switch. \citet{KambourovManovskii2009} demonstrate that occupational tenure returns dominate firm and industry returns (12--20\% for five years), establishing that human capital is occupation-specific. \citet{Pedulla2016} documents overqualification stigma: employers penalize overqualified applicants with callback penalties approximately as severe as one year of unemployment, creating genuine downward barriers.

A critical methodological issue is selection bias from observing only completed transitions \citep{Heckman1979}. Workers blocked by credential requirements do not appear in transition data. Our institutional barrier estimates reflect friction experienced by successful transitioners, not blocking probability for all who attempt. This selection likely causes us to underestimate true credential effects.


\subsection{Shock Characterization}
\label{subsec:lit-shocks}


The RBTC paradigm operationalizes shocks through the routine task intensity index: $\text{RTI}_k = \ln(R_k) - \ln(M_k) - \ln(A_k)$ \citep{AutorDorn2013}. This discrete classification captures first-order patterns but obscures within-category heterogeneity and graded spillovers.

AI exposure measures provide occupation-level indices of technology impact. \citet{FeltenRajSeamans2021} construct AIOE by linking AI capabilities to occupational abilities. \citet{EloundouEtAl2024} use GPT-4 ratings of task-level exposure. \citet{Webb2020} matches patent text to job descriptions, providing precedent for patent-to-task exposure constructions as an alternative to O*NET-only measures. These measures identify which occupations contain exposed tasks; they do not characterize \emph{where} exposed workers can go. Recent work includes the ILO's refined global occupational exposure index \citep{ILO2025} and the White House Council of Economic Advisers' empirical comparison of AI exposure measures \citep{WhiteHouse2024}.

A gap remains in systematic methodology for tracking capability frontier dynamics---how technology capabilities evolve and which task domains they enter. Current approaches use point-in-time snapshots rather than dynamic capability tracking.


\subsection{Adjustment Mechanisms}
\label{subsec:lit-mechanisms}


Switching cost estimation spans an order of magnitude in US data. \citet{LeeWolpin2006} find sectoral switching costs of 0.50--0.75$\times$ annual wages using Census/NLSY data, noting that occupational switching within sectors is less costly than sectoral switching. \citet{ArtucChaudhuriMcLaren2010} estimate 4.7--6.5$\times$ annual wages using CPS sectoral flows. \citet{ArtucMcLaren2015} find 0.99--1.54$\times$ with job quality controls. \citet{DixCarneiro2014} estimates 1.4--2.7$\times$ using Brazilian administrative data with full identification of worker fixed effects.

\citet{Traiberman2019} represents the structural frontier. Using Danish linked employer-employee data, she estimates switching costs averaging five years of income (IQR: two years) with task-dimension-specific prices. Her methodology requires individual-level wages, complete work histories, and full population coverage---data unavailable in US household surveys. A critical finding: ignoring human capital accumulation biases switching cost estimates upward by a factor of three.

Gravity models of labor flows provide benchmarks for predictive accuracy. \citet{CortesGallipoli2018} use US CPS aggregate flows with DOT task measures, finding task-related costs explain no more than 15\% of total transition costs across most occupation pairs. Our gravity model partial $R^2$ of 2.55\% (centroid) is consistent with task distance playing a similarly modest role. \citet{SiminiEtAl2021} achieve correlations of 0.3--0.5 for geographic mobility predictions. These benchmarks contextualize our supply-demand decomposition.


\subsection{Feasibility Versus Realization}
\label{subsec:lit-feasibility}


The distinction between where workers \emph{can} go and where they \emph{do} go has theoretical precedent. \citet{BarseghyanEtAl2021} distinguish the feasible set (all theoretically available alternatives) from the choice set (what agents actually consider and select) in discrete choice modeling. \citet{YiMuellerStegmaier2017} examine skill transferability as structural determinant of reallocation potential. \citet{LeeWolpin2006} explicitly assess the relative importance of labor demand and supply factors for sectoral reallocation, finding demand-side factors drove service-sector growth.

Policy applications legitimize feasibility-focused analysis. Workforce development programs (WIOA) target training to structurally feasible occupations. Career pathway systems map feasible ladders based on skill adjacency. Green transition planning uses skill distance to assess worker reallocation from fossil fuel to renewable sectors. Feasibility is the supply-side input that, combined with demand-side factors, determines equilibrium outcomes.


% ==================================================================
% @AGENT: STABLE THEORY — LOW CHANGE FREQUENCY
% Definitions are referenced by NAME (e.g., "Effective Distance Decomposition")
% not by number (numbers shift). Do not modify without explicit instruction.
% Key objects: Task domain, Occupation measures, d_eff decomposition, Wasserstein
% ==================================================================

\section{Framework Architecture and Module Specifications}
\label{sec:theory}


This section defines the framework architecture and specifies each module. The structure is: (i) the framework as organizing architecture; (ii) module specifications for $\mathcal{T}$, $\mathcal{I}$, $\mathcal{S}$, and $\mathcal{M}$; (iii) operational definitions for empirical implementation.


\subsection{The Task-Space Framework}
\label{subsec:framework}


\begin{definition}[Task-Space Framework]
\label{def:framework}
The \emph{task-space framework} $\mathcal{F}: (\mathcal{T}, \mathcal{I}, \mathcal{S}, \mathcal{M}) \to (\Delta\boldsymbol{\rho}, \Delta \mathbf{L}, \Delta \mathbf{W})$ is the research program target: a complete structural mapping from task-space representation ($\mathcal{T}$), institutions ($\mathcal{I}$), technology shocks ($\mathcal{S}$), and equilibrium mechanisms ($\mathcal{M}$) to occupation-specific changes in task distributions, employment, and wages.
\end{definition}

Here $\Delta\boldsymbol{\rho}$ denotes occupation-specific changes in task distributions (how the bundle of tasks each occupation performs evolves), while $\Delta \mathbf{L}$ and $\Delta \mathbf{W}$ denote occupation-specific changes in employment and wages. These are the observable outcomes that policy analysis targets.

This paper validates candidate specifications for $\mathcal{T}$ and $\mathcal{I}$, integrates external $\mathcal{S}$, and identifies requirements for $\mathcal{M}$. The framework is an organizing architecture, not an estimation target. Its value is taxonomic: it clarifies which components a complete model would require, allowing us to test components in isolation and identify which combinations survive falsifiable checks.


\subsection{$\mathcal{T}$ Module: Task Representation}
\label{subsec:t-module}


The $\mathcal{T}$ module encodes occupation skill content as a metric space. The core question: what makes two occupations ``close'' in terms of worker mobility?


\begin{definition}[Task domain]
\label{def:task-domain}
The \emph{task domain} is a compact metric-measure space $\mathcal{T} = (T, d, \mu)$, where $T$ is a set of tasks, $d: T \times T \to \mathbb{R}_{\geq 0}$ is a metric encoding similarity, and $\mu$ is a reference measure for economic mass.
\end{definition}

Distance is the primitive notion of similarity: tasks $x, y$ are close if capabilities and complements transfer at low friction. The framework uses this geometry to encode a prior on spillovers: shocks affect nearby tasks more strongly.


To build intuition before the formal specification of occupations as measures, we develop two toy examples.


\begin{example}[The Task Line]
\label{ex:task-line}
Consider a simple task domain $T = [0, 1]$ representing a single dimension of complexity (e.g., ``manual'' to ``cognitive''). The metric $d(x, y) = |x - y|$ captures the cost of retooling from task $x$ to task $y$. The reference measure $\mu$ is uniform, representing an equal density of economic value across the complexity spectrum.
\end{example}


\begin{remark}[Manifold vs Graph]
\label{rem:graph-vs-manifold}
The manifold representation assumes continuous similarity: tasks $x, y$ at small distance $d(x,y)$ are economically interchangeable at low friction. An alternative representation treats activities as discrete nodes in a graph, where adjacency (shared occupation usage) rather than continuous distance determines relatedness. The empirical program tests which representation better characterizes economic outcomes.
\end{remark}


\begin{definition}[Occupations as measures]
\label{def:occupations}
An occupation $j$ is represented by a probability measure $\rho_j$ over $T$ with density $f_j(x) = d\rho_j / d\mu(x)$, interpreted as the intensity with which occupation $j$ uses task $x$.
\end{definition}


\begin{example}[Occupations on the Line (toy illustration)]
\label{ex:occupations-line}
To build intuition, consider two occupations on the task line $T = [0,1]$. ``Occupation A'' (Specialist) is concentrated at the manual end, modeled as a Gaussian measure centered at $x=0.2$ with narrow variance $\sigma_A^2 = 0.01$. ``Occupation B'' (Generalist) is centered at $x=0.8$ but has wider variance $\sigma_B^2 = 0.04$, so its support extends partially into the middle range. The overlap integral $\int \rho_A(x) \rho_B(x)\, d\mu(x)$ quantifies their task-space relatedness. This is intuition only; the empirical implementation uses 2,087 activities and 894 occupations.
\end{example}


\begin{assumption}[Admissible task representation]
\label{assump:admissible-representation}
An empirical representation of the task domain is admissible if, for at least one validation target not used in construction (e.g., wage comovement, worker mobility), occupation-pair relatedness in activity space predicts stronger observed relatedness.
\end{assumption}


\subsubsection{Candidate Geometries}


The 2$\times$2 methodology comparison isolates two design choices: (i) how to represent inter-task similarity (embedding-based vs.\ O*NET-based), and (ii) how to aggregate task-level information to occupation distance (distributional vs.\ centroid). We specify the candidates:

\textbf{Embedding ground metric} encodes semantic task substitutability: tasks with similar meaning in natural language are treated as economically substitutable. Cosine distance in sentence embedding space (MPNet, 768 dimensions) captures that ``operating forklift'' $\approx$ ``driving delivery vehicle'' even when distinctly labeled.

\textbf{Identity ground metric} treats all tasks as equally different: $d(a, b) = \mathbf{1}[a \neq b]$. This is the implicit assumption of non-embedding structured-vector measures that treat tasks as independent dimensions.

\textbf{Wasserstein (distributional) aggregation} measures minimum transformation cost:
\begin{definition}[Wasserstein distance between occupations]
\label{def:wasserstein}
\[
W(\rho_i, \rho_j) = \min_{\gamma \in \Pi(\rho_i, \rho_j)} \sum_{a,b} \gamma(a,b) \cdot d(a,b),
\]
where $\gamma$ is a transport plan in $\Pi(\rho_i, \rho_j)$ (joint distributions with marginals $\rho_i$ and $\rho_j$), and $d(a,b)$ is the ground distance between tasks.
\end{definition}

The transport plan $\gamma(a,b)$ represents optimal skill transformation---the allocation of retraining investment across task pairs. Total Wasserstein distance is total transformation cost. This ``earth mover's'' interpretation is economically intuitive: workers reshaping one skill distribution into another.

\textbf{Centroid (non-distributional) aggregation} computes occupation embedding centroids and measures cosine distance between them:
\[
d_{\mathrm{centroid}}(i, j) = 1 - \cos\left(\bar{e}_i, \bar{e}_j\right), \quad \text{where } \bar{e}_i = \frac{\sum_a \rho_i(a) \cdot e_a}{\sum_a \rho_i(a)}
\]
and $e_a$ is the embedding vector for task $a$.


\begin{remark}[Embedding vs.\ distributional: which drives improvement?]
\label{rem:embedding-vs-distributional}
The 2$\times$2 design isolates these effects. Results (\cref{subsec:geometry-comparison}) show:
\begin{itemize}
    \item \textbf{Embedding effect:} 74.9\% improvement in pseudo-$R^2$ (embedding vs.\ O*NET methods)
    \item \textbf{Distributional effect:} Wasserstein adds nothing beyond centroid; centroid marginally outperforms
\end{itemize}
The embedding ground metric---semantic task substitutability---is the mechanism. Wasserstein provides theoretical grounding (minimum transformation cost) but no additional empirical advantage after correcting for diagonal bias in the distance matrix; simpler centroid methods approximate the geometry closely ($\rho = 0.95$ correlation in distance rankings) and marginally outperform.
\end{remark}


\subsection{$\mathcal{I}$ Module: Institutional Structure}
\label{subsec:i-module}


The $\mathcal{I}$ module encodes non-skill barriers to mobility. The central theoretical contribution is the decomposition of effective distance into semantic and institutional components.


\begin{definition}[Effective distance decomposition]
\label{def:effective-distance}
Define the \emph{effective distance} between occupations $i$ and $j$ as:
\[
d_{\mathrm{eff}}(i,j) = d_{\mathrm{sem}}(i,j) + \lambda \cdot d_{\mathrm{inst}}(i,j),
\]
where $d_{\mathrm{sem}}(i,j) \geq 0$ is semantic distance (embedding-based Wasserstein or centroid), $d_{\mathrm{inst}}(i,j) \geq 0$ is institutional distance (job zone difference plus certification barrier), and $\lambda \geq 0$ is the relative weight.
\end{definition}


\begin{remark}[The brain surgeon / butcher problem]
\label{rem:surgeon-butcher}
Brain surgeon and butcher both involve cutting tissue with precision tools. Semantically, they may appear similar. Economically, they are vastly distant: a butcher cannot become a brain surgeon without years of training and credentialing. Semantic distance may be moderate; institutional distance is maximal. The decomposition correctly identifies them as economically distant despite superficial task similarity.
\end{remark}


\begin{remark}[Semantic operationalization affects decomposition]
\label{rem:operationalization-decomposition}
The choice of semantic distance measure affects correlation with institutional distance:
\begin{itemize}
    \item Non-embedding baselines: low correlation with institutional distance (largely orthogonal)
    \item Embedding-based measures: $\rho(d_{\mathrm{wasserstein}}, d_{\mathrm{inst}}) = 0.249$ (moderate correlation)
\end{itemize}
Embedding-based distances capture skill transformation costs that correlate with training requirements. The residual institutional component after controlling for embedding distance may capture genuinely non-skill barriers: credentialing gates, licensing, discrimination, network effects. However, this interpretation requires caution; see \cref{subsec:cps-mobility} for discussion.
\end{remark}


\begin{definition}[Asymmetric institutional distance]
\label{def:asymmetric-distance}
The symmetric $d_{\mathrm{inst}}$ can be decomposed into directional components:
\[
d_{\mathrm{inst}}^{\mathrm{up}}(i,j) = \max(0, \mathrm{Zone}_j - \mathrm{Zone}_i) + \gamma \cdot \max(0, \mathrm{Cert}_j - \mathrm{Cert}_i),
\]
\[
d_{\mathrm{inst}}^{\mathrm{down}}(i,j) = \max(0, \mathrm{Zone}_i - \mathrm{Zone}_j) + \gamma \cdot \max(0, \mathrm{Cert}_i - \mathrm{Cert}_j).
\]
By construction, $d_{\mathrm{inst}}^{\mathrm{up}}(i,j) + d_{\mathrm{inst}}^{\mathrm{down}}(i,j) = d_{\mathrm{inst}}(i,j)$.
\end{definition}


\begin{remark}[Theoretical predictions for asymmetry]
\label{rem:asymmetry-prediction}
\textbf{Credential-gate theory} predicts $\beta_{\mathrm{up}} \gg \beta_{\mathrm{down}}$: licensing restricts entry but not exit \citep{Jackson2023}. \textbf{Human capital specificity theory} predicts $\beta_{\mathrm{up}} \approx \beta_{\mathrm{down}}$: occupation-specific capital is lost in any move \citep{KambourovManovskii2009,GathmannSchonberg2010}. \textbf{Overqualification stigma} \citep{Pedulla2016} creates genuine downward barriers. \Cref{subsubsec:asymmetric-test} tests these predictions.
\end{remark}


\begin{remark}[Barriers as friction versus gates]
\label{rem:friction-vs-gates}
A critical distinction: ``friction'' is continuous and affects all transitions proportionally; ``gates'' are discontinuous and block certain transitions entirely. Our institutional distance coefficient is estimated from completed transitions only. Workers blocked by credential requirements do not appear in our data. This selection bias means $\gamma_{\mathrm{inst}}$ estimates friction experienced by successful transitioners, not the prohibitive effect on blocked workers. The low ratio $\gamma_{\mathrm{inst}}/\gamma_{\mathrm{sem}} = 0.019$ (\cref{subsec:cps-mobility}) likely understates true credential effects.
\end{remark}


\subsection{$\mathcal{S}$ Module: Shock Profiles}
\label{subsec:s-module}


The $\mathcal{S}$ module characterizes which occupations contain tasks exposed to technological change. A shock profile specifies intensity across the task domain.


\begin{definition}[Shock profile]
\label{def:shock-profile}
A shock profile $I: T \to \mathbb{R}$ specifies innovation intensity at each task. Occupation-level exposure aggregates via the task distribution:
\[
E_j = \int_T I(x)\, d\rho_j(x).
\]
\end{definition}


Two approaches to shock profile construction:
\begin{enumerate}
    \item \textbf{External anchoring:} Use independently constructed measures (AIOE, Eloundou GPT-4 ratings) as $I$. This paper adopts external anchoring for $\mathcal{S}$.
    \item \textbf{Endogenous estimation:} Estimate $I$ from observed outcome variation. Requires identifying assumptions about shock exogeneity.
\end{enumerate}

The key empirical question: do external shock profiles integrate coherently with geometry-based feasibility? If AIOE identifies exposed occupations and embedding-based distance identifies compatible destinations, the constructs should be distinct (low correlation) but combinable (joint model improves probabilistic scoring).


\subsection{$\mathcal{M}$ Module: Adjustment Mechanisms}
\label{subsec:m-module}


The $\mathcal{M}$ module specifies how shocks propagate to outcomes. A complete $\mathcal{M}$ would include: switching costs (investment required to change occupations), search frictions (time and matching efficiency), wage adjustment (equilibrium price response), and general equilibrium spillovers (how one occupation's shock affects others).


\begin{definition}[Switching cost specification]
\label{def:switching-cost}
Following \citet{Traiberman2019}, switching costs can be parameterized as:
\[
C(o, o') = \exp\left[\Gamma^M + \sum_i (\upsilon^{o'}_i - \upsilon^o_i) \cdot \Gamma_i\right],
\]
where $\upsilon^o$ is occupation $o$'s task intensity vector and $\Gamma_i$ prices movement along each task dimension. This accommodates our decomposition by adding credential indicators as additional cost dimensions.
\end{definition}


Identification of $\mathcal{M}$ parameters requires individual-level wage data at the moment of transition. With aggregate occupation-mean wages, we cannot distinguish entry wages from incumbent wages, causing misidentification (\cref{subsec:cost-calibration}). Structural estimation following \citet{Traiberman2019} requires administrative data unavailable for the US.

This paper adopts \textbf{external calibration} for $\mathcal{M}$: anchoring switching cost magnitudes to structural estimates from \citet{DixCarneiro2014}. This is an interim approach. Full $\mathcal{M}$ development requires richer data (LEHD, administrative records) and is future work.

\begin{remark}[Supply-demand separation in the framework architecture]
\label{rem:supply-demand-separation}
The framework's modular structure separates supply-side geometry ($\mathcal{T}$, $\mathcal{I}$ modules) from demand-side mechanisms ($\mathcal{M}$ module). Geometry defines the feasible choice set---which transitions are possible given skill requirements. Demand determines selection within that set---which destinations actually receive inflows. The supply-demand decomposition (\cref{subsec:pathway-validation}) provides empirical grounding for this architectural separation: demand dominates aggregate flows ($\rho \approx 0.80$), while geometry captures pathway direction ($\rho \approx 0.12$).
\end{remark}


\subsection{Structural Grounding}
\label{subsec:structural-grounding}


The decomposition $d_{\mathrm{eff}} = d_{\mathrm{sem}} + \lambda \cdot d_{\mathrm{inst}}$ finds structural support in \citet{Traiberman2019}, who parameterizes switching costs as linear in task-space distance with dimension-specific prices. Estimated switching costs average five years of income (IQR: two years), with intrasectoral occupation switching 28\% costlier than pure sectoral switching.

\begin{remark}[Relation to other structural models]
\label{rem:structural-comparison}
\citet{ArtucChaudhuriMcLaren2010} parameterize sector-switching costs at 6.5 times annual wages but treat all transitions equivalently. \citet{KambourovManovskii2009} demonstrate occupational tenure returns dominate firm and industry returns but assume binary specificity rather than continuous task distance. \citet{Lazear2009} provides the skill-weights approach: all skills are general but occupations weight them differently.
\end{remark}


\subsection{Propagation of Shocks Across Tasks}
\label{subsec:propagation}


A technology shock to one activity may spill over to similar activities. We model this as distance-weighted averaging.


\begin{definition}[Propagation mechanisms]
\label{def:propagation-mechanisms}
\textbf{(A) Kernel propagation:} 
$[\widetilde{\mathcal{K}}_d f](x) = \frac{\int_T k(d(x,y)) f(y)\, d\mu(y)}{\int_T k(d(x,y))\, d\mu(y)}$, where $k(\cdot)$ is positive and decreasing in distance.

\textbf{(B) Identity propagation:} No spillover: $[\mathcal{I} f](x) = f(x)$.
\end{definition}


\subsection{Aggregation and Outcome Equations}
\label{subsec:aggregation}


\begin{definition}[Occupation-level exposure functionals]
\label{def:exposure-functionals}
Given technological state $(A_t, C_t)$ and occupation measures $\{\rho_j\}$:
\[
E^A_{j,t} \equiv \int_T A_t(x)\, d\rho_j(x), \qquad E^C_{j,t} \equiv \int_T \log C_t(x)\, d\rho_j(x).
\]
The pipeline is: shock profile on tasks $\to$ propagation across tasks $\to$ aggregation to occupations.
\end{definition}


\begin{definition}[Reduced-form outcome equations]
\label{def:reduced-form}
Let $w_{j,t}$ and $L_{j,t}$ denote wages and employment:
\[
\Delta \log w_{j,t} = \beta_w E^A_{j,t} + \gamma_w E^C_{j,t} + \bm{X}'_{j,t} \delta_w + u_{j,t},
\]
\[
\Delta \log L_{j,t} = \beta_L E^A_{j,t} + \gamma_L E^C_{j,t} + \bm{X}'_{j,t} \delta_L + v_{j,t}.
\]
Coefficients are conditional associations for comparing exposure measures.
\end{definition}


\section{Data Infrastructure and Module Operationalization}
\label{sec:data}


This section operationalizes the framework modules using O*NET data. We describe: (i) model space axes; (ii) discrete approximation of the task manifold; (iii) distance matrix construction; (iv) external data integration for $\mathcal{S}$ and $\mathcal{M}$ modules.


\subsection{Model Space Axes}
\label{subsec:model-space}


The paper explores three axes:
\begin{enumerate}
    \item \textbf{Representation axis:} How to represent inter-task similarity (embedding-based [MPNet, MiniLM] vs.\ O*NET-based [importance weights, binary indicators])
    \item \textbf{Aggregation axis:} How to aggregate task-level information (distributional [Wasserstein] vs.\ centroid [cosine on centroids])
    \item \textbf{Shock axis:} How to characterize technology (external AIOE, generic capability, modality-specific)
\end{enumerate}
Each empirical test constrains this space. The 2$\times$2 methodology comparison (\cref{subsec:geometry-comparison}) establishes that representation choice (embedding vs.\ O*NET) dominates aggregation method (Wasserstein vs.\ centroid). Validated combinations advance; rejected combinations are pruned.


\subsection{Discrete Approximation of the Task Manifold}
\label{subsec:discrete-approx}


We use O*NET Detailed Work Activities (DWAs) as the task unit. O*NET provides 2,087 DWAs linked to 894 occupations via task statements.


\begin{definition}[Empirical activity domain]
\label{def:empirical-activity-domain}
Let $T_n = \{a_1, \ldots, a_n\}$ denote the set of DWAs ($n = 2{,}087$). The reference measure $\mu$ is uniform. The metric $d$ is constructed from text embeddings: cosine distance in MPNet embedding space (768 dimensions).
\end{definition}


\begin{remark}[Constructing occupation measures from O*NET]
\label{rem:rho-construction}
The occupation measure $\rho_j$ is constructed from O*NET ratings \citep{onet2024dwa,onet2024scales}:

\textbf{(A) Raw weight:} For DWAs, importance is derived via linked tasks: $w_j(a) = \max_{t \in \mathrm{Tasks}(j,a)} \mathrm{IM}_{j,t}$.

\textbf{(B) Filtering:} Set $w_j(a) = 0$ if O*NET flags ``Recommend Suppress'' or ``Not Relevant.''

\textbf{(C) Normalization:} $\rho_j(a) = w_j(a) / \sum_{a'} w_j(a')$.

\textbf{(D) Coverage:} 894 occupations with valid measures. Sparsity: 99.1\% of occupation-activity weights are zero. Median activities per occupation: 19.
\end{remark}


\subsection{Distance Matrix Construction}
\label{subsec:distance-construction}


\textbf{Embedding-based Wasserstein distance:}
\begin{enumerate}
    \item Embed 2,087 DWAs using MPNet (all-mpnet-base-v2, 768 dimensions)
    \item Compute cosine distance matrix $D_{\mathrm{act}}$ over activities (ground metric)
    \item Row-normalize occupation-activity matrix to probability measures
    \item Compute Wasserstein distance for each occupation pair using POT library
\end{enumerate}
Computation time: 84 seconds for full $894 \times 894$ matrix.

\textbf{Embedding-based centroid distance:}
\begin{enumerate}
    \item Embed 2,087 DWAs using MPNet
    \item Compute importance-weighted centroid for each occupation: $\bar{e}_j = \sum_a \rho_j(a) \cdot e_a$
    \item Compute cosine distance between occupation centroids
\end{enumerate}

\textbf{O*NET-based cosine distance:}
\begin{enumerate}
    \item Represent each occupation as vector of DWA importance weights
    \item Compute cosine distance directly on importance vectors (no embedding)
\end{enumerate}

\textbf{O*NET-based Euclidean distance:}
\begin{enumerate}
    \item Represent each occupation as vector of DWA importance weights
    \item Compute Euclidean distance on importance vectors
\end{enumerate}

\textbf{Institutional distance:}
\[
d_{\mathrm{inst}}(i,j) = |\mathrm{Zone}_i - \mathrm{Zone}_j| + \gamma \cdot |\mathrm{Cert}_i - \mathrm{Cert}_j|
\]
Job Zone (1--5) has 100\% coverage; certification importance (1--5) has 61.1\% coverage with median imputation.


\subsection{External Data Integration}
\label{subsec:external-data}


\textbf{AI Occupational Exposure (AIOE):} \citet{FeltenRajSeamans2021} scores linked via SOC crosswalk. Coverage: 93.5\% of occupations in CPS sample.

\textbf{OES Wages:} Bureau of Labor Statistics Occupational Employment and Wage Statistics, 2023. Coverage: 98.2\% of sample occupations. Mean annual wage: \$71,992.

\textbf{BLS Openings:} Bureau of Labor Statistics Occupational Outlook Handbook projected annual openings. Coverage: 94.3\% of sample occupations. Used for demand decomposition analysis.

\textbf{CPS Transitions:} IPUMS CPS Basic Monthly Files, 2015-01 through 2024-09 (excluding 2020--2021 for COVID disruption). Full sample: $n = 89{,}329$ (after persistence and O*NET mapping filters). 2024 subsample ($n = 8{,}880$) used for out-of-period model comparison.


\begin{table}[h!]
\centering
\caption{Design Axis Specifications}
\label{tab:design-axes}
\begin{tabular}{lll}
\toprule
\textbf{Axis} & \textbf{Primary} & \textbf{Comparison/Robustness} \\
\midrule
Task representation & Embedding (MPNet) & O*NET importance weights \\
Occupation distance & \textbf{Cosine on centroids} & Wasserstein on embeddings, O*NET-based \\
Ground metric & Embedding similarity & Identity (all tasks equally different) \\
Shock profile & External (AIOE) & Generic capability \\
\bottomrule
\end{tabular}
\end{table}


% ==================================================================
% @AGENT: LOCKED RESULTS
% These findings are verified. Cite as facts. See LEDGER.md for structured lookup.
% Modification requires new experimental evidence, not reinterpretation.
% ==================================================================

\section{Module Validation and Model Space Constraints}
\label{sec:validation}


This section reports module validation results. Each subsection follows the structure: objective, specification, results, module status. Tests are organized as checkpoints for framework component validation.


\subsection{$\mathcal{T}$ Module Validation: 2$\times$2 Methodology Comparison}
\label{subsec:geometry-comparison}


\textbf{Objective:} Isolate the contribution of embedding-based task representation from the contribution of distributional (Wasserstein) aggregation. Test which design choice drives improvement in probabilistic scoring of realized transitions.

\textbf{Design:} A 2$\times$2 factorial comparison crossing:
\begin{itemize}
    \item \textbf{Representation:} Embedding-based (MPNet sentence embeddings) vs.\ O*NET-based (importance weights)
    \item \textbf{Aggregation:} Distributional (Wasserstein optimal transport) vs.\ non-distributional (cosine/Euclidean on vectors or centroids)
\end{itemize}

\textbf{Specification:} Conditional logit for destination choice using $n = 89{,}329$ CPS occupation transitions:
\[
U_{kj} = \alpha \cdot (-d_{\mathrm{sem}}(i,j)) + \beta \cdot (-d_{\mathrm{inst}}(i,j)) + \varepsilon_{kj}
\]
The only difference between specifications is the semantic distance measure.

Following \citet{McFadden1978}, we estimate using sampled alternatives: for each observed transition, the choice set consists of the realized destination plus 10 randomly drawn alternatives ($J = 11$). Pseudo-$R^2$ is computed against this sampled choice set; values are not directly comparable to full-choice-set estimates over all 447 destinations.

\medskip
\noindent\textit{Identification note.} This reduced-form choice model tests whether a distance metric assigns higher probability to realized transitions than alternatives. It does not separately identify worker retraining costs, employer screening preferences, vacancy distributions, or unobserved institutional constraints. The model provides a probabilistic scoring comparison between distance constructions; mechanism interpretations are stated as ``consistent with,'' not ``identified as.''

Two additional estimation caveats apply. First, pseudo-$R^2$ values in Tables 2--5 are computed in-sample: the same 89,329 transitions used for estimation are used for performance comparison. With two estimated parameters ($\alpha$, $\beta$) across 89,329 choice occasions, overfitting risk is minimal, and the 2024 out-of-period comparison for the shock integration module (\cref{subsec:shock-integration}) provides out-of-sample confirmation. However, the reported pseudo-$R^2$ values should be interpreted as relative performance across specifications rather than out-of-sample prediction accuracy. Second, when multiple O*NET-SOC codes map to a single Census occupation code, task distributions and distances are aggregated using unweighted means. Employment-weighted aggregation would better represent the typical worker's task profile within each Census code but requires linking OES employment data to the crosswalk, which we leave to future work. Third, the sampled alternative set does not exclude the origin occupation. In approximately 2\% of choice sets, the origin appears as an alternative with zero distance. A robustness check excluding the origin from the alternative set yields $\alpha = 7.704$ (vs.\ 7.394), a change of 4.2\%, confirming this design choice does not materially affect results.


\begin{table}[h!]
\centering
\caption{2$\times$2 Distance Metric Comparison: Embedding vs.\ O*NET $\times$ Aggregation Method}
\label{tab:2x2-comparison}
\begin{tabular}{llcccc}
\toprule
\textbf{Representation} & \textbf{Aggregation} & $\alpha$ & Pseudo-$R^2$ & $\Delta$ vs baseline \\
\midrule
Embedding (MPNet) & Cosine centroid & 7.40 & 14.08\% & +132\% \\
Embedding (MPNet) & Wasserstein & 8.39 & 13.8\% & +128\% \\
O*NET importance & Cosine & 4.55 & 8.05\% & +33\% \\
O*NET importance & Euclidean & 9.76 & 6.06\% & baseline \\
\bottomrule
\end{tabular}
\end{table}


\textbf{Results:} The 2$\times$2 comparison isolates two effects:
\begin{itemize}
    \item \textbf{Embedding effect:} Embedding-based methods (13.8--14.1\% pseudo-$R^2$) vs.\ O*NET-based methods (6.1--8.1\%) represents a 74.9\% relative improvement. This is the dominant effect.
    \item \textbf{Distributional effect:} Wasserstein (13.8\%) vs.\ cosine-on-centroids (14.08\%)---Wasserstein adds nothing beyond centroid; centroid marginally outperforms. This strengthens the paper's core argument that the embedding representation, not the distributional treatment, drives the gain.
\end{itemize}

The embedding ground metric---semantic task substitutability---drives the improvement. Workers transitioning between occupations behave as if semantically similar tasks are substitutable, even when distinctly labeled in O*NET. The distributional treatment provides theoretical grounding (minimum transformation cost interpretation) but limited additional empirical advantage.


\subsubsection{Ground Metric Validation}


To further isolate the embedding contribution, we compare Wasserstein distance computed with the embedding ground metric versus an identity ground metric (all tasks equally different):


\begin{table}[h!]
\centering
\caption{Ground Metric Validation: Semantic Similarity vs.\ Identity}
\label{tab:ground-metric}
\begin{tabular}{lcccc}
\toprule
\textbf{Ground Metric} & $\alpha$ & $\gamma$ & Pseudo-$R^2$ & $\Delta$ \\
\midrule
Identity (all tasks equally different) & 4.71 & 0.32 & 7.52\% & baseline \\
Embedding (semantic similarity) & 8.39 & 0.154 & 13.76\% & +83\% \\
\bottomrule
\end{tabular}
\end{table}


The embedding ground metric improves explanatory power by over 80\% (+83\%). Knowing that semantically similar tasks (e.g., ``operating forklift'' $\approx$ ``driving delivery vehicle'') are substitutable adds substantial predictive value beyond raw task overlap.


\subsubsection{Distance Correlation Analysis}


\begin{table}[h!]
\centering
\caption{Distance Matrix Correlations (Spearman $\rho$)}
\label{tab:distance-correlations}
\begin{tabular}{lc}
\toprule
\textbf{Pair} & \textbf{Spearman $\rho$} \\
\midrule
Wasserstein (embedding) vs.\ Cosine centroid (embedding) & 0.95 \\
Wasserstein (embedding) vs.\ Wasserstein (identity) & 0.36 \\
Wasserstein (embedding) vs.\ Euclidean (O*NET) & 0.28 \\
Cosine centroid vs.\ Wasserstein (identity) & 0.35 \\
\bottomrule
\end{tabular}
\end{table}


\textbf{Key finding:} Wasserstein and cosine-on-centroids (both embedding-based) produce nearly identical occupation distance rankings ($\rho = 0.95$). The embedding ground metric fundamentally reshapes the geometry: correlation with identity-based Wasserstein is only 0.36. Centroid averaging inherits the embedding structure, not the task overlap structure.


\subsubsection{Mechanism: Semantic Task Substitutability}


The embedding ground metric encodes \emph{semantic task substitutability}: the insight that tasks with similar meaning in natural language are economically interchangeable. This explains why:
\begin{itemize}
    \item ``Operating forklift'' and ``driving delivery vehicle'' should be treated as similar (vehicle operation skills transfer)
    \item ``Analyzing financial data'' and ``preparing budget reports'' should be treated as similar (quantitative analysis skills transfer)
    \item Traditional O*NET measures that treat these as independent dimensions miss the transferability
\end{itemize}

The sentence embedding model (MPNet) learns this semantic structure from a large text corpus. When applied to O*NET task descriptions, it identifies which tasks involve similar capabilities even when labeled differently. Workers' revealed mobility patterns validate this semantic similarity as economically relevant.


\textbf{Module status: $\mathcal{T}$ VALIDATED for embedding-informed distance.} The embedding ground metric (semantic task substitutability) is the validated mechanism. Wasserstein optimal transport provides theoretical grounding as minimum transformation cost but does not outperform simpler centroid methods after correcting for diagonal bias; centroid-based distance marginally outperforms. For practical applications where computational simplicity matters, cosine distance on embedding centroids is preferred.


\subsection{$\mathcal{I}$ Module Validation: Institutional Distance}
\label{subsec:cps-mobility}


\textbf{Objective:} Test whether institutional distance adds predictive power conditional on semantic distance. This validates the separability assumption of the $\mathcal{I}$ module.

\textbf{Data:} $n = 89{,}329$ verified CPS occupation transitions (2015--2019, 2022--2024). Persistence filter applied to remove measurement error. Our estimand is mobility patterns among verified/sustained transitions, not attempted transitions.


\begin{table}[h!]
\centering
\caption{CPS Sample Construction Pipeline}
\label{tab:cps-pipeline}
\begin{tabular}{lrrp{5cm}}
\toprule
\textbf{Stage} & \textbf{Records} & \textbf{Retention} & \textbf{Notes} \\
\midrule
Raw CPS extract & 10,072,119 & 100\% & 83 monthly samples \\
Employment filter & 4,407,432 & 43.8\% & Age 18--65, employed \\
Consecutive month pairs & 2,353,725 & 53.4\% & CPSIDP linked \\
Raw transitions & 157,606 & 6.74\% & OCC$(t) \neq$ OCC$(t-1)$ \\
Persistence filter & 106,116 & 67.3\% & Verified transitions \\
O*NET mapping & 89,329 & 84.2\% & Census $\to$ O*NET \\
\bottomrule
\end{tabular}
\end{table}


\begin{table}[h!]
\centering
\caption{Conditional Logit Results: Embedding-Centroid (Primary)}
\label{tab:clogit-primary}
\begin{tabular}{lcccc}
\toprule
\textbf{Parameter} & \textbf{Coefficient} & \textbf{Std. Error} & \textbf{$t$-statistic} & \textbf{$p$-value} \\
\midrule
$\alpha$ (semantic) & 7.404 & 0.036 & 204.0 & $< 10^{-100}$ \\
$\beta$ (institutional) & 0.139 & 0.0041 & 33.7 & $< 10^{-100}$ \\
\bottomrule
\end{tabular}
\end{table}


\textbf{Results:} Both coefficients are positive and highly significant. Transitions are more frequently observed to destinations with lower semantic distance (skill transformation cost) and lower institutional distance (credential barriers). The decomposition assigns higher probability to realized transitions.

\textbf{Module status: $\mathcal{I}$ VALIDATED.} Institutional distance is a significant predictor conditional on semantic distance. Separability holds.


\subsubsection{Directional Asymmetry Test}
\label{subsubsec:asymmetric-test}


\textbf{Objective:} Test whether institutional barriers are asymmetric---whether credentials function as one-way gates.


\begin{table}[h!]
\centering
\caption{Asymmetric Conditional Logit: Geometry Comparison}
\label{tab:asymmetric-geometry}
\begin{tabular}{lccccc}
\toprule
\textbf{Geometry} & $\alpha$ (semantic) & $\beta_{\mathrm{up}}$ & $\beta_{\mathrm{down}}$ & Ratio & LR $\chi^2$ \\
\midrule
O*NET cosine & 5.691 & 0.282 & 0.270 & 1.04 & 4.33 \\
Embedding Wasserstein & 9.001 & 0.171 & 0.081 & 2.11 & 242.1 \\
\bottomrule
\end{tabular}
\end{table}


\begin{table}[h!]
\centering
\caption{Robustness of Asymmetry Ratio Under Embedding Wasserstein}
\label{tab:asymmetric-robustness}
\begin{tabular}{llcc}
\toprule
\textbf{Variant} & \textbf{Description} & \textbf{Ratio} & \textbf{95\% CI} \\
\midrule
S1 & Baseline & 2.11 & [1.80, 2.42] \\
S2 & Job Zone only & 1.36 & [1.23, 1.50] \\
R1 & Thick cells ($\geq 5$ transitions) & 2.79 & [2.33, 3.26] \\
R2 & Prime age (25--54) & 1.12 & [0.96, 1.29] \\
R3 & Excluding outlier occupations & 0.06 & [0.01, 0.11] \\
\bottomrule
\end{tabular}
\end{table}


\textbf{Results:} Directional asymmetry is sample-dependent. Under O*NET cosine, barriers appear symmetric (ratio = 1.04). Under embedding Wasserstein, baseline shows asymmetry (ratio = 2.11, $\chi^2 = 242$), but prime-age workers show symmetric barriers (ratio = 1.12), and excluding outlier occupations eliminates asymmetry (ratio = 0.06).

\textbf{Interpretation:} Neither pure credential-gate nor pure symmetric-friction theories are universally supported. This heterogeneity is an informative constraint: $\mathcal{I}$ module specifications should accommodate sample-dependent patterns. Three mechanisms may explain prime-age symmetry: overqualification discrimination \citep{Pedulla2016}, human capital specificity \citep{KambourovManovskii2009,GathmannSchonberg2010}, and Job Zones measuring distance rather than gates \citep{KleinerXu2020}.


\subsection{$\mathcal{S}$ Module Validation: Shock Profile Linkage}
\label{subsec:shock-integration}


\textbf{Objective:} Test whether external AI exposure scores integrate coherently with geometry-based feasibility. If AIOE identifies exposed occupations and embedding distance identifies compatible destinations, these should be distinct constructs.

\textbf{Sample:} Full sample $n = 89{,}329$; 2024 subsample $n = 8{,}880$ used for out-of-period model comparison.


\begin{table}[h!]
\centering
\caption{Shock Integration Validation}
\label{tab:shock-integration}
\begin{tabular}{lcc}
\toprule
\textbf{Metric} & \textbf{Value} & \textbf{Interpretation} \\
\midrule
AIOE coverage & 93.5\% & High linkage success \\
AIOE--centroid correlation & $-$0.016 & Empirically distinct \\
\midrule
\multicolumn{3}{l}{\textit{Model comparison (2024 out-of-period transitions):}} \\
Geometry log-likelihood & $-41{,}745$ & --- \\
Historical baseline & $-65{,}625$ & --- \\
$\Delta$LL (Geometry vs Historical) & $+23{,}879$ & Strong improvement \\
\bottomrule
\end{tabular}
\end{table}


\textbf{Results:} AIOE and embedding distance are orthogonal ($r = -0.02$), confirming they capture distinct constructs. Shock profiles identify which occupations contain exposed tasks; geometry identifies \emph{where} exposed workers can feasibly transition. Geometry substantially outperforms historical baselines in probabilistic scoring ($\Delta \text{LL} = +23{,}879$).

\textbf{Module status: $\mathcal{S}$ VALIDATED.} External shock profiles and feasibility geometry combine coherently. The orthogonality confirms complementary constructs; out-of-period comparison confirms predictive validity.


\subsection{$\mathcal{M}$ Module Probe: Switching Cost Calibration}
\label{subsec:cost-calibration}


\textbf{Objective:} Attempt endogenous switching cost identification; establish external calibration as interim approach when endogenous fails.


\subsubsection{Endogenous Estimation (Failed)}


\begin{table}[h!]
\centering
\caption{Endogenous Wage Coefficient Estimation}
\label{tab:endogenous-wage}
\begin{tabular}{lcccc}
\toprule
\textbf{Parameter} & \textbf{M1 (log)} & \textbf{M2 (level)} & \textbf{M3 (rank)} \\
\midrule
$\gamma_{\mathrm{sem}}$ & 8.935 & 8.877 & 8.935 \\
$\gamma_{\mathrm{inst}}$ & 0.136 & 0.148 & 0.136 \\
$\beta_{\mathrm{wage}}$ & $-0.063$ & 0.011 & $-0.063$ \\
$t$-stat (wage) & $-6.7$ & 1.2 & $-6.7$ \\
\bottomrule
\end{tabular}
\end{table}


\textbf{Diagnosis:} Negative $\beta_{\mathrm{wage}}$ in the log specification indicates model misidentification. Workers appear to \emph{avoid} high-wage destinations after controlling for skill distance. The problem: occupation-mean wages do not capture the entry wages switchers actually receive. Incumbents earn tenure premiums; switchers start lower. Using mean wages as the wage signal introduces systematic bias.

\textbf{Data requirement identified:} Individual-level wages at transition (LEHD, administrative records) are required for endogenous $\mathcal{M}$ estimation.


\subsubsection{External Calibration (Adopted)}


\begin{table}[h!]
\centering
\caption{External Switching Cost Calibration}
\label{tab:external-calibration}
\begin{tabular}{lcc}
\toprule
\textbf{Parameter} & \textbf{Value} & \textbf{Source} \\
\midrule
Calibration anchor & 2.0 wage-years & Dix-Carneiro (2014) median \\
Median centroid distance & 0.229 & Training sample \\
\textbf{SC per unit centroid} & \textbf{8.74 wage-years} & Derived \\
SC per unit centroid & \$628,581 & Using OES 2023 mean wage \\
\bottomrule
\end{tabular}
\end{table}


\textbf{Method:} Anchor to \citet{DixCarneiro2014}'s median estimate (2.0 wage-years for typical sectoral transition). With median centroid distance 0.229, this implies 8.74 wage-years per unit centroid distance; equivalently, with median Wasserstein distance 0.521, 3.84 wage-years per unit Wasserstein distance. US estimates span 0.5$\times$ to 6.5$\times$ annual wages; ordinal rankings of transition difficulty are invariant across calibration sources.


\begin{table}[h!]
\centering
\caption{Sensitivity Analysis: Calibration Source}
\label{tab:calibration-sensitivity}
\begin{tabular}{lcc}
\toprule
\textbf{Source} & \textbf{Anchor (wage-years)} & \textbf{SC/unit Wasserstein} \\
\midrule
Lee \& Wolpin (2006) & 0.75 & $\sim$1.4 \\
Dix-Carneiro mid (adopted) & 2.0 & 3.84 \\
Dix-Carneiro upper & 2.7 & $\sim$5.2 \\
Artuc et al.\ (2010) & 6.5 & $\sim$12.5 \\
\bottomrule
\end{tabular}
\end{table}


\textbf{Caveat:} External calibration requires maintained assumptions: (i) cross-country applicability (Brazilian $\to$ US); (ii) sectoral-to-occupational translation. US estimates span 0.5$\times$ to 6.5$\times$ annual wages---an order of magnitude. Results should be interpreted with this uncertainty. Ordinal rankings of transition difficulty are invariant across the calibration range; absolute magnitudes scale proportionally.


\subsubsection{Example Transitions (Smell Test)}


\begin{table}[h!]
\centering
\caption{Example Transition Costs}
\label{tab:example-transitions}
\begin{tabular}{llccc}
\toprule
\textbf{From} & \textbf{To} & \textbf{$d_{\mathrm{wass}}$} & \textbf{Cost (wage-yrs)} & \textbf{Interpretation} \\
\midrule
Cashiers & Retail Salespersons & 0.35 & 1.35 & Adjacent roles \\
Software Developers & Data Scientists & 0.44 & 1.69 & Related technical \\
Accountants & Software Developers & 0.55 & 2.12 & Cross-domain \\
Cashiers & Registered Nurses & 0.65 & 2.49 & Credential-requiring \\
\bottomrule
\end{tabular}
\end{table}


\textbf{Interpretation:} Costs scale appropriately with intuitive skill distance. Adjacent roles (cashier $\to$ retail sales) are cheap; credential-requiring transitions (cashier $\to$ nurse) are expensive; cross-domain technical transitions fall in between. However, these costs are conditional on successful realized transition; workers blocked by credential requirements (estimated at $\sim$24\% by \citealt{Jackson2023}) do not appear in our data. The smell test passes for completed transitions.

\textbf{Module status: $\mathcal{M}$ CALIBRATED VIA EXTERNAL BENCHMARK.} Endogenous estimation requires richer data. Calibrated costs provide planning numbers for policy analysis.


\subsection{Scope Validation: Supply-Demand Decomposition}
\label{subsec:pathway-validation}


\textbf{Objective:} Quantify the supply-demand decomposition of occupational transitions. Test whether geometry captures supply-side feasibility while demand determines realized flows.

\textbf{Methodology:} We distinguish two metrics:
\begin{itemize}
    \item \textbf{Per-origin pathway direction:} For each origin occupation, correlate model-predicted destination rankings (based on inverse centroid distance) with observed destination rankings, then average across origins.
    \item \textbf{Aggregate inflow prediction:} Correlate predicted total inflows by destination with observed total inflows.
\end{itemize}


\begin{table}[h!]
\centering
\caption{Supply-Demand Decomposition of Transition Patterns}
\label{tab:demand-decomposition}
\begin{tabular}{lcc}
\toprule
\textbf{Predictor} & \textbf{Spearman $\rho$} & \textbf{Interpretation} \\
\midrule
Demand-only (BLS openings) & 0.798 & Dominant for aggregate inflows \\
Geometry-only (1/distance) & $\approx$0.00 & Negligible for aggregate \\
Per-origin geometry & 0.118 & Modest for pathway direction \\
\bottomrule
\end{tabular}
\end{table}


\subsubsection{Performance Battery}
\label{subsubsec:performance-battery}


Beyond log-likelihood, we report ranking adequacy and dispersion diagnostics to characterize what the geometry captures.


\begin{table}[h!]
\centering
\caption{Performance Battery: Ranking Adequacy and Dispersion}
\label{tab:performance-battery}
\begin{tabular}{lcc}
\toprule
\textbf{Metric} & \textbf{Value} & \textbf{Interpretation} \\
\midrule
Mean Percentile Rank (MPR) & 0.74 (median 0.83) & Top 26\% on average \\
Realized Cumulative Mass (RCM) & 0.56 & 56\% mass at/above realized \\
Effective consideration set ($N_{\mathrm{eff}}$) & 207 of 447 (46\%) & Diffuse distribution \\
\midrule
Top-5 overlap & Not reported & $N_{\mathrm{eff}}/J = 0.46 > 0.3$ \\
Top-10 overlap & Not reported & Too diffuse for top-$k$ \\
\bottomrule
\end{tabular}
\end{table}


\textbf{Interpretation:} Mean Percentile Rank (MPR) = 0.74 indicates that realized destinations rank in the top 26\% of predicted probability on average; half rank in the top 17\% (median 0.83). Realized Cumulative Mass (RCM) = 0.56 means 56\% of probability mass is assigned to destinations at or above the realized choice. The effective consideration set ($N_{\mathrm{eff}} = 207$ of 447 destinations, 46\%) indicates the model distributes probability across nearly half the destination space.

Top-$k$ overlap is not reported because with $N_{\mathrm{eff}}/J = 0.46$, distributions are too diffuse for top-$k$ overlap to be meaningful. This explains why Top-5 overlap = 0 is expected under a feasibility framing rather than a point-prediction framing.

These metrics support a feasibility interpretation: the model identifies a broad consideration set of skill-compatible destinations ($N_{\mathrm{eff}} \approx 207$) and ranks the realized destination highly within that set (MPR = 0.74). It does not predict a single destination---nor should it, given that demand factors dominate realized flows ($\rho = 0.80$).


\begin{table}[h!]
\centering
\caption{Pathway Identification: 2024 Out-of-Period Validation}
\label{tab:pathway-validation}
\begin{tabular}{lcc}
\toprule
\textbf{Metric} & \textbf{Value} & \textbf{Interpretation} \\
\midrule
N observed transitions & 8,880 & 2024 subsample \\
N origins evaluated & 233 & Full sample \\
\textbf{Per-origin Spearman $\rho$} & \textbf{0.118} & Modest pathway direction signal \\
KL divergence & 0.854 & Moderate divergence \\
\bottomrule
\end{tabular}
\end{table}


\textbf{Interpretation:} The decomposition reveals that aggregate occupational inflows are overwhelmingly demand-determined ($\rho = 0.80$). Geometry-based feasibility provides modest but statistically significant information about pathway direction at the origin level ($\rho \approx 0.12$), consistent with \citet{CortesGallipoli2018}'s finding that task-related costs account for no more than 15\% of total transition costs. The framework architecture's separation of supply-side geometry ($\mathcal{T}$, $\mathcal{I}$ modules) from demand-side mechanisms ($\mathcal{M}$ module) is consistent with this pattern: geometry defines the feasible choice set, while demand determines selection within that set. This parallels the distinction between consideration sets and revealed choice in discrete choice models \citep{BarseghyanEtAl2021} and supply-demand decompositions in structural labor economics \citep{LeeWolpin2006}.


\begin{table}[h!]
\centering
\caption{Observed vs.\ Predicted Top-5 Destinations (AI-Exposed Origins)}
\label{tab:observed-vs-predicted}
\begin{tabular}{cl|cl}
\toprule
\textbf{Rank} & \textbf{Observed} & \textbf{Rank} & \textbf{Predicted} \\
\midrule
1 & Managers, All Other & 1 & Secondary School Teachers \\
2 & Executive Secretaries & 2 & Preschool Teachers \\
3 & Office Clerks, General & 3 & Special Ed Teachers \\
4 & Elementary Teachers & 4 & Social Service Managers \\
5 & Bookkeeping Clerks & 5 & Advertising Managers \\
\bottomrule
\end{tabular}
\end{table}


\textbf{Interpretation:} Predicted destinations are skill-compatible but credential-gated or capacity-constrained. The ratio $\gamma_{\mathrm{inst}}/\gamma_{\mathrm{sem}} = 0.019$ means credentials barely register---because blocked transitions are unobserved. This is selection bias inherent in revealed preference estimation (\cref{rem:friction-vs-gates}), not model failure.

\textbf{Scope validated:} The geometry correctly captures supply-side feasibility constraints. Combined with the demand decomposition ($\rho \approx 0.80$ for aggregate flows), this quantifies the scope claim: the framework measures where workers \emph{can} go given skill architecture, not where they \emph{do} go given labor market conditions.


\subsection{Aggregate Flow Analysis: Gravity Model}
\label{subsec:gravity}


\textbf{Objective:} Test whether task distance explains aggregate bilateral occupation flows, and compare to the \citet{CortesGallipoli2018} benchmark.

\textbf{Specification:} Standard gravity model on occupation pairs:
\[
\ln(\text{Flow}_{ij} + 1) = \alpha + \beta_1 \ln(\text{Emp}_i) + \beta_2 \ln(\text{Emp}_j) + \beta_3 \cdot d(i,j) + \varepsilon_{ij}
\]
Sample: 199,362 occupation pairs (including zeros). Mass-only $R^2 = 22.06\%$.


\begin{table}[h!]
\centering
\caption{Gravity Model: Distance Metric Comparison}
\label{tab:gravity-comparison}
\begin{tabular}{lccc}
\toprule
\textbf{Distance Metric} & $\beta_{\mathrm{distance}}$ & $t$-statistic & Partial $R^2$ \\
\midrule
Embedding $\times$ Centroid & $-0.577$ & $-82.2$ & 2.55\% \\
Embedding $\times$ Wasserstein & $-1.041$ & $-96.2$ & 3.46\% \\
O*NET $\times$ Cosine & $-1.556$ & $-92.4$ & 3.20\% \\
O*NET $\times$ Euclidean & $-0.514$ & $-25.3$ & 0.25\% \\
\bottomrule
\end{tabular}
\end{table}


\textbf{Results:} Task distance explains a modest share of aggregate bilateral flow variation. The Embedding $\times$ Centroid specification achieves partial $R^2 = 2.55\%$; Embedding $\times$ Wasserstein achieves 3.46\%. All distance coefficients are negative and highly significant: greater task distance reduces bilateral flows.

\textbf{Comparison to Cortes-Gallipoli:} Our partial $R^2$ values (0.25\%--3.46\%, with Embedding $\times$ Centroid at 2.55\%) are consistent with their finding that task-specific costs account for no more than 15\% of total transition costs across most occupation pairs. The dominant explanatory power in gravity models comes from origin and destination characteristics (captured by employment mass terms), not pairwise task distance.

\textbf{Ranking divergence from conditional logit:} Distance metrics rank differently in gravity than conditional logit. Notably, cosine similarity on sparse O*NET vectors---which performs worst for individual choice prediction---is competitive for aggregate flows (partial $R^2 = 3.20\%$ vs.\ 3.46\% for Wasserstein). This pattern reflects extensive vs.\ intensive margin dynamics.


\begin{remark}[Extensive vs.\ intensive margin interpretation]
\label{rem:extensive-intensive}
The divergence between conditional logit and gravity rankings is informative about occupational mobility mechanisms:
\begin{itemize}
    \item \textbf{Conditional logit} conditions on workers who switched---it measures the intensive margin (which destination among connected occupations?). Individual workers navigating occupational space use continuous task similarity gradients that embedding-based distances capture.
    \item \textbf{Gravity} includes zero-flow pairs---it captures both extensive margin (are these occupations connected at all?) and intensive margin. The near-binary structure of sparse O*NET cosine (78\% of pairs at maximum distance) captures whether occupation pairs exchange any workers, which dominates aggregate flow variation.
\end{itemize}
Embedding-based distances excel at fine-grained choice among connected occupations. Simpler measures capturing ``connected vs.\ not'' structure are competitive for aggregate flows.
\end{remark}


\textbf{Interpretation:} Individual workers navigating occupational space use continuous task similarity gradients that embedding-based distances capture. Aggregate flow patterns depend primarily on network connectivity structure---whether occupations exchange any workers---where simpler measures suffice. Both findings are consistent with task distance playing a modest but significant role in occupational mobility, and both align with \citet{CortesGallipoli2018}'s structural estimates.


\subsection{Complementary Validations}
\label{subsec:complementary}


\subsubsection{Wage Comovement}


\begin{table}[h!]
\centering
\caption{Wage Comovement Validation}
\label{tab:wage-comovement}
\begin{tabular}{lccc}
\toprule
\textbf{Measure} & $\beta$ & $t$ & $R^2$ \\
\midrule
Normalized embedding overlap & 0.337 & 7.39 & 0.00523 \\
Binary Jaccard & 0.471 & 8.00 & 0.00167 \\
\bottomrule
\end{tabular}
\end{table}


Semantic overlap is associated with occupation-pair wage correlations ($R^2 = 0.52\%$), outperforming binary Jaccard by 3$\times$. The signal is statistically robust but quantitatively modest.


\subsubsection{RTI Construct Validity}


\begin{table}[h!]
\centering
\caption{Correlation: Semantic Exposure vs.\ RTI Components}
\label{tab:rti-correlation}
\begin{tabular}{lcc}
\toprule
\textbf{Component} & \textbf{Pearson $r$} & \textbf{$p$-value} \\
\midrule
RTI (composite) & $-0.060$ & 0.303 \\
Routine & $+0.002$ & 0.978 \\
Abstract & $+0.088$ & 0.130 \\
Manual & $+0.090$ & 0.124 \\
\bottomrule
\end{tabular}
\end{table}


Embedding-based semantic exposure is essentially orthogonal to Autor-Dorn RTI. The geometry captures \emph{mobility friction}---where can workers realistically transition?---rather than \emph{automation susceptibility}---which tasks technology will displace. These are distinct constructs. The marginal automation prediction result ($\Delta R^2 = 2.2\%$ over RTI, $p = 0.075$) reflects this distinction.


\subsection{Model Space Status Summary}
\label{subsec:model-space-status}


\begin{table}[h!]
\centering
\caption{Module Status Summary}
\label{tab:module-status}
\begin{tabular}{llll}
\toprule
\textbf{Module} & \textbf{Status} & \textbf{Validated} & \textbf{Remaining} \\
\midrule
$\mathcal{T}$ & Operational & Embedding $\gg$ O*NET (74.9\%) & Domain-specific embeddings \\
$\mathcal{I}$ & Operational & Separable, significant & Licensing vs.\ job zones; gate modeling \\
$\mathcal{S}$ & Validated & AIOE linkage; $\Delta$LL = +23,879 & Modality-specific profiles \\
$\mathcal{M}$ & Preliminary & Costs calibrated externally & Endogenous estimation; GE extension \\
Scope & Quantified & Supply-demand decomposition & Demand integration (Phase 0.8) \\
\bottomrule
\end{tabular}
\end{table}


\subsection{Structural Stability: Pre/Post COVID Comparison}
\label{subsec:covid-stability}


\textbf{Objective:} Test whether embedding-based task distance captures structural skill transferability or contingent revealed preferences. COVID provides a natural experiment: the largest labor market disruption in decades. If task-distance penalties reflect structural skill constraints, they should be invariant across this disruption; if they reflect contingent labor market conditions, they should shift.

\textbf{Specification:} We split the CPS sample into pre-COVID (2015--2019, $n = 60{,}225$) and post-COVID (2022--2024, $n = 29{,}104$) periods, excluding 2020--2021 to avoid acute disruption effects. We estimate the same conditional logit specification used for $\mathcal{T}$ module validation (embedding centroid) separately by period, then conduct a likelihood ratio test for structural break.


\begin{table}[h!]
\centering
\caption{Structural Stability: Pre/Post COVID Comparison (Centroid Specification)}
\label{tab:covid-stability}
\begin{tabular}{lcccccc}
\toprule
\textbf{Period} & $\alpha$ & \textbf{SE} & $\gamma$ & \textbf{SE} & \textbf{Pseudo-$R^2$} \\
\midrule
Pre-COVID (2015--2019) & 7.394 & 0.044 & 0.146 & 0.005 & 14.1\% \\
Post-COVID (2022--2024) & 7.358 & 0.063 & 0.144 & 0.007 & 13.9\% \\
\bottomrule
\end{tabular}
\end{table}


\textbf{Aggregate stability:} The likelihood ratio test for structural break yields $\chi^2(2) = 0.67$, $p = 0.72$---we cannot reject parameter equality across periods. The task-distance penalty ($\alpha$) changes by less than 1\% (7.394 $\to$ 7.358). This invariance across the largest labor market disruption in decades supports the interpretation that embedding-based distance captures structural skill transferability constraints rather than contingent revealed preferences. This finding is consistent with \citet{PizzinelliShibata2023}'s evidence that aggregate occupational mismatch remained stable through COVID despite massive employment disruption.


\subsubsection{Heterogeneity by Teleworkability}


Despite aggregate stability, compositional heterogeneity exists. We interact task distance with destination teleworkability \citep{DingelNeiman2020} and period:
\[
U_{kj} = \alpha \cdot (-d_{\mathrm{sem}}) + \gamma \cdot (-d_{\mathrm{inst}}) + \delta_3 \cdot (-d_{\mathrm{sem}} \times \mathrm{Telework}_j) + \delta_4 \cdot (-d_{\mathrm{sem}} \times \mathrm{Telework}_j \times \mathrm{Post}) + \varepsilon_{kj}
\]

\textbf{Results:} $\delta_3 = -0.561$ ($p < 0.0001$): Teleworkable destinations exhibit stronger task-distance penalties in both periods---transitions to remote-capable occupations are more skill-constrained. $\delta_4 = -0.086$ ($p = 0.010$): This penalty increased post-COVID. The finding is robust to excluding 2022 (immediate post-pandemic); qualitatively similar results obtain using cosine-on-centroids rather than Wasserstein distance.

\textbf{Interpretation:} We interpret this heterogeneity through \citet{ModestinoShoagBallance2020}'s ``upskilling'' framework. Remote work expansion dramatically increased applicant pools for teleworkable positions \citep{HsuTambe2022}, enabling employers to be more selective. The elevated task-distance penalty reflects heightened hiring standards---employers screening more stringently on skill match when applicant pools are larger---rather than changes in underlying skill requirements. This interpretation is consistent with \citet{EmanuelHarrington2024}'s evidence on selection effects in remote work. We use ``elevated hiring standards'' rather than ``tightened geometry'' to emphasize that this is an equilibrium response to supply conditions, not a change in structural skill constraints.

\textbf{Caveat:} We cannot fully distinguish employer selectivity from genuine skill requirement updates. If remote work changed which tasks occupations actually perform (e.g., more written communication, less in-person coordination), the elevated penalty could reflect real changes in task content. The aggregate stability finding ($\Delta\alpha < 1\%$) suggests composition effects rather than structural changes, but this remains an interpretation, not an identification.

\textbf{Implication:} The aggregate stability supports structural interpretation of embedding-based distance---task-similarity constraints are invariant across major labor market disruption. The heterogeneity demonstrates that the framework can detect economically meaningful variation, positioning it as a tool for substantive labor market analysis rather than measurement infrastructure alone.


% ==================================================================
% @AGENT: FRONTIER — ACTIVE RESEARCH
% Contains pending tests and future work. Check LEDGER.md for current sprint.
% This section changes frequently as paths are completed or deferred.
% ==================================================================

\section{Extensions and Deferred Investigations}
\label{sec:extensions}


The modular architecture enables incremental extension. This section records completed extensions, identified but unexecuted investigations, and deferred future work. The research program continues---validated modules provide foundations for systematic exploration.


\textbf{Completed extensions:}
\begin{itemize}
    \item 2$\times$2 methodology comparison (\cref{subsec:geometry-comparison}): Embedding effect (74.9\%) dominates distributional effect (centroid marginally outperforms Wasserstein)
    \item Ground metric validation (\cref{subsec:geometry-comparison}): Embedding ground metric +83\% over identity
    \item Asymmetric institutional distance (\cref{subsubsec:asymmetric-test}): Heterogeneous, sample-dependent asymmetry
    \item RTI construct validity (\cref{subsec:complementary}): Orthogonal to RTI ($r = -0.06$)
    \item Shock integration (\cref{subsec:shock-integration}): AIOE linkage validated
    \item Switching cost calibration (\cref{subsec:cost-calibration}): External benchmark applied
    \item Supply-demand decomposition (\cref{subsec:pathway-validation}): Demand $\rho = 0.80$, geometry $\rho = 0.12$
    \item Gravity model analysis (\cref{subsec:gravity}): Centroid partial $R^2 = 2.55\%$, consistent with Cortes-Gallipoli
\end{itemize}

\textbf{Identified extensions (specified, not executed):}
\begin{itemize}
    \item Domain-specific embedding comparison (JobBERT vs.\ MPNet)
    \item Licensing as alternative barrier measure
    \item Prospective AI evaluation with candidate shock profiles
    \item Institutional absorption interaction test (effect $\times$ licensing intensity)
\end{itemize}

\textbf{Deferred extensions (future phases):}
\begin{itemize}
    \item Full retrospective battery (Autor-Dorn replication)
    \item Modality-specific shock profiles
    \item Task composition outcome validation
    \item Structural estimation (Traiberman-style)
\end{itemize}


\subsection{Dynamic Accumulation and Reversals}
\label{subsec:dynamic}


The baseline construction is static. For settings where shock accumulation or reversal matters, define:
\[
A_t = \delta A_{t-1} + I_t + \varepsilon_t,
\]
where $\delta \in [0,1]$ governs persistence. This extension enables modeling of cumulative technology exposure and potential skill depreciation.


\subsection{Asymmetric Institutional Distance (Completed)}
\label{subsec:asymmetric}


The asymmetric extension has been tested (\cref{subsubsec:asymmetric-test}). The finding is geometry-sensitive and sample-dependent: under O*NET cosine, barriers appear symmetric (ratio $= 1.04$); under embedding Wasserstein, the baseline shows asymmetry (ratio $= 2.11$), but this asymmetry does not survive all robustness checks.

This finding discriminates between theoretical models in a nuanced way. Neither pure credential-gate theory ($\beta_{\mathrm{up}} \gg \beta_{\mathrm{down}}$) nor pure human capital specificity theory ($\beta_{\mathrm{up}} \approx \beta_{\mathrm{down}}$) is universally supported. The evidence is consistent with heterogeneous barriers: the pattern varies by worker age (prime-age workers show symmetric barriers) and occupation type (excluding outlier occupations eliminates asymmetry).

Three mechanisms may explain the heterogeneity:
\begin{enumerate}
    \item \textbf{Occupation-specific human capital} \citep{KambourovManovskii2009}: 12--20\% wage premiums forfeited in any switch
    \item \textbf{Overqualification stigma} \citep{Pedulla2016}: Employers penalize overqualified applicants
    \item \textbf{Job Zones measure preparation, not permission} \citep{KleinerXu2020}: Licensing affects mobility symmetrically
\end{enumerate}

The finding constrains future framework specifications: geometry choice affects institutional inference, and directional barrier patterns should be tested for sample-dependence.


\subsection{Licensing as Alternative Barrier Measure}
\label{subsec:licensing-extension}


Job Zones and occupational licensing measure fundamentally different constructs. Job Zones capture human capital investment requirements---years of education, training time, accumulated experience. Licensing creates legal permission requirements---government authorization to work in an occupation.

\citet{Jackson2023} demonstrates asymmetry using CPS licensing data: licensing reduces entry probability with no significant effect on exit probability. This is the one-way gate pattern that Job Zones fail to capture. \citet{KleinerXu2020} find different results with symmetric effects on both entry and exit---the literature remains unsettled on whether licensing creates one-way or two-way barriers.

Future work should test the decomposition using: (i) CPS licensing supplement (2015+) matched to Institute for Justice policy data; (ii) state-level licensing variation following \citet{KleinerXu2020}; (iii) specific licensed occupations with binding credential requirements (nursing compact variation, bar exam pass rates, medical licensing pathways). These specifications should reveal whether the asymmetry that simple credentialing theory predicts appears with more targeted institutional measures.


\subsection{Augmentation Channel}
\label{subsec:augmentation}


The framework carries $C_t$ (complementarity field) and $E^C$ (augmentation exposure) for symmetry. Estimation requires task-level productivity complementarities not available in O*NET. This channel is deferred pending data development.


\subsection{Supervised Dimension Extraction}
\label{subsec:supervised}


An alternative to unsupervised text embeddings is supervised extraction of economically-relevant dimensions. Given theoretical priors about dimensions that matter (routine/non-routine, cognitive/manual), one can train linear probes to identify concept directions and project all activities onto these dimensions. This approach trades the generality of unsupervised embeddings for targeted economic interpretability.


\subsection{Worker Mobility as Extended Validation}
\label{subsec:mobility-extension}


The CPS mobility test uses destination \emph{choice}. An extension would model transition \emph{rates}: do occupation pairs with lower $d_{\mathrm{eff}}$ have higher observed transition probabilities? This tests whether the geometry is associated with the \emph{likelihood} of switching, not just the direction conditional on switching.


\subsection{Historical Regime Signatures}


To summarize historical episodes without forcing a single narrative, \cref{tab:regimes} records regime signatures as hypotheses about shock support, displacement vs augmentation, and primary mechanism.


\begin{table}[h!]
\centering
\begin{tabular}{|l|l|l|l|l|}
\hline
\textbf{Regime} & \textbf{Shock support ($I_t$)} & \textbf{$A$ vs.\ $C$} & \textbf{Mechanism} & \textbf{Wage wedge} \\
\hline
Electrification (1900--1940) & Manual/craft tasks & Mixed & TBD & Context-dependent \\
\hline
Postwar scaling (1945--1970) & Broad/diffuse & Augmentation & TBD & High pass-through \\
\hline
ICT/RBTC (1980--2007) & Routine-cognitive core & Displacement & Both informative & Declining pass-through \\
\hline
AI era (2020--) & Analytic/generative & TBD & TBD & TBD \\
\hline
\end{tabular}
\caption{Historical regimes as parameter signatures.}
\label{tab:regimes}
\end{table}


\subsection{Prospective Evaluation: Generative AI}
\label{subsec:prospective}


The $\mathcal{S}$ module requires shock profile candidates for future validation. The prospective evaluation applies the framework to generative AI. The object is comparison among pre-specified shock profiles, not identification of a single ``true'' profile.


\subsubsection{Candidate Shock Profiles}


Each candidate defines $I^{(v)}(a)$ for activities $a \in T_n$.

\medskip
\noindent\textbf{$v_1$ (capability-only):}
\begin{itemize}
    \item Construction: Score each activity on ``language/cognitive capability relevance'' using a rule-based mapping.
    \item Positive loading on activities tagged as involving information processing, written comprehension, oral expression; negative loading on physical/manual activities.
\end{itemize}

\noindent\textbf{$v_2$ (capability $+$ structure):}
\begin{itemize}
    \item Construction: $v_1$ loadings plus additional positive weight on activities tagged as structured/repetitive.
    \item ``Structured'' operationalized via O*NET ``Structured vs Unstructured'' rating.
    \item Hypothesis: structure amplifies AI exposure beyond raw capability.
\end{itemize}

\noindent\textbf{$v_3$ (external anchor---PRIMARY):}
\begin{itemize}
    \item Construction: Derive from an independent source not using O*NET geometry.
    \item Options: (a) Human expert ratings; (b) Published rubric mapping (Eloundou et al.\ 2024); (c) Patent-activity matching following \citet{Webb2020}.
    \item Designated as primary profile ex ante.
\end{itemize}


\subsubsection{Profile Selection Protocol}


\noindent\textit{Primary-profile approach (baseline):}
\begin{itemize}
    \item Designate $v_3$ as primary profile ex ante.
    \item $v_1$ and $v_2$ are robustness checks and hypothesis tests.
    \item No profile selection based on outcome fit.
\end{itemize}


\subsubsection{Implications for AI-Based Exposure Measures}


The CPS mobility findings have direct implications for how AI shock profiles should be deployed:

\begin{enumerate}
    \item \textbf{Any task-space framework must incorporate institutional distance.} Semantic similarity alone overstates the feasibility of worker reallocation. A model that assumes workers displaced by AI will flow to semantically similar occupations without accounting for credentialing barriers will systematically mischaracterize adjustment paths.
    
    \item \textbf{AI exposure profiles should distinguish ``can do'' from ``allowed to do.''} If generative AI displaces tasks in high-credential occupations (e.g., legal research, medical documentation), workers may face institutional barriers to moving into other occupations even if their semantic skills transfer.
    
    \item \textbf{Counterfactual reallocation simulations should respect $d_{\mathrm{eff}}$.} When simulating how workers will adjust to AI shocks, the framework should weight destination occupations by $\exp(-d_{\mathrm{eff}})$, not just $\exp(-d_{\mathrm{sem}})$.
    
    \item \textbf{Supply-side feasibility is necessary but not sufficient.} The demand decomposition ($\rho \approx 0.80$ for demand, $\rho \approx 0.12$ for geometry) demonstrates that realized transitions depend heavily on labor market conditions. AI impact projections should incorporate both feasibility constraints and demand-side factors.
\end{enumerate}


\subsection{Empirical Paths Forward}
\label{subsec:paths-forward}


The research program follows pre-committed decision rules. Each path specifies a test, expected outcomes, and implications for model space.


\subsubsection{Path A: 2$\times$2 Methodology Comparison --- RESOLVED}


\textbf{Objective:} Isolate embedding contribution from distributional (Wasserstein) contribution.

\textbf{Status: RESOLVED.} Embedding effect (74.9\%) dominates distributional effect (centroid marginally outperforms Wasserstein). The embedding ground metric (semantic task substitutability) is the validated mechanism.

\textbf{Pre-committed interpretation applied:}
\begin{table}[h!]
\centering
\begin{tabular}{lll}
\toprule
\textbf{Outcome} & \textbf{Pre-committed Interpretation} & \textbf{Actual Result} \\
\midrule
Embedding $\gg$ O*NET & Semantic task substitutability matters & \textbf{Confirmed (74.9\%)} \\
Wasserstein $\gg$ Centroid & Distributional treatment adds value & Not confirmed (centroid marginally outperforms) \\
\bottomrule
\end{tabular}
\end{table}

\textbf{Implication:} Embedding-informed distance is adopted as primary $\mathcal{T}$ specification. Wasserstein provides theoretical grounding; simpler centroid methods approximate the geometry closely ($\rho = 0.95$) and marginally outperform.


\subsubsection{Path B: Domain-Specific Embedding Comparison --- DEFERRED}


\textbf{Objective:} Test whether domain-specific embeddings (JobBERT-v2, trained on ESCO skills) outperform general-purpose embeddings (MPNet).

\textbf{Status: DEFERRED.} With embedding-informed distance validated as primary geometry, domain-specific comparison is lower priority. The embedding effect (74.9\%) is large; further embedding optimization unlikely to produce comparable improvement.

\textbf{Pre-committed interpretations:}
\begin{table}[h!]
\centering
\begin{tabular}{lll}
\toprule
\textbf{Outcome} & \textbf{Interpretation} & \textbf{Next Step} \\
\midrule
JobBERT $\gg$ MPNet & Domain-specific training captures economic structure & Adopt JobBERT \\
MPNet $\gg$ JobBERT & General semantic structure sufficient & Stay with MPNet \\
JobBERT $\approx$ MPNet & Embedding choice doesn't matter for this task & Prefer MPNet (standard) \\
\bottomrule
\end{tabular}
\end{table}


\subsubsection{Path C: RTI Construct Validity --- COMPLETED}


\textbf{Status: COMPLETED.} Embedding-based semantic exposure was correlated with Autor-Dorn RTI at occ1990dd level. Result: $r = -0.060$ (not significant).

\textbf{Implication:} The semantic geometry captures mobility friction, not automation susceptibility---a construct economically orthogonal to classic RBTC routine exposure. This explains the marginal automation prediction result and positions the framework as complementary to, rather than redundant with, task-based automation measures.


\subsubsection{Path D: Modality-Specific Shocks --- DEFERRED}


\textbf{Objective:} Construct capability profiles distinguishing code generation, reasoning, agentic workflows. Map to O*NET tasks.

\textbf{Status: DEFERRED.} Comprehensive implementation requires 6--9 months (taxonomy design, task annotation, benchmark mapping). One-month pilot (2--3 occupations, $\sim$200 tasks) could establish methodology feasibility.


\subsubsection{Path E: Task Composition Outcomes --- DEFERRED}


\textbf{Objective:} Replace employment with skill mention changes in job postings \citep{HampoleEtAl2025}.

\textbf{Data requirement:} Lightcast/Revelio access (institutional license required).

\textbf{Status: DEFERRED.} Data not currently available. Documented as research program extension.


\subsubsection{Path F: Asymmetric Test with Embedding Wasserstein --- COMPLETED}


\textbf{Status: COMPLETED.} The asymmetric barriers test was run using embedding Wasserstein geometry. Baseline ratio = 2.11 (vs.\ 1.04 under O*NET cosine). Robustness analysis reveals sample-dependent heterogeneity: prime-age workers show ratio = 1.12; excluding outlier occupations yields ratio = 0.06. See \cref{subsubsec:asymmetric-test} for full results.

\textbf{Implication:} Geometry choice affects institutional inference. Directional asymmetry is heterogeneous, not universal.


\subsubsection{Path G: Supply-Demand Decomposition --- COMPLETED}


\textbf{Status: COMPLETED.} Demand-only $\rho = 0.798$ for aggregate inflows; geometry-only $\rho \approx 0.00$ for aggregate; per-origin geometry $\rho = 0.118$ for pathway direction.

\textbf{Implication:} Aggregate flows are demand-dominated. Geometry captures supply-side constraints (necessary condition for transition), not demand-side outcomes (sufficient condition). The framework architecture's separation of feasibility from realization is consistent with observed patterns.


\subsubsection{Path H: Retrospective Battery --- PARTIALLY COMPLETE}


\textbf{Objective:} Test embedding geometry against historical employment polarization patterns in the Autor-Dorn setting (1980--2005).

\textbf{Status: PARTIALLY COMPLETE.} Test B (CZ-level polarization) executed. Tests A and C deferred due to data/methodology constraints. See \cref{app:retrospective} for full results.

\textbf{Key result:} Embedding geometry adds modest incremental explanatory power ($\Delta R^2 = 5.5\%$ for operator employment decline), but the orthogonal component of CSH does not systematically predict employment polarization (1+, 0$-$, 4(0) across five outcomes).

\textbf{Implication:} The weak incremental explanatory power is expected given construct distinction. CSH captures task similarity/mobility friction; RTI captures automation susceptibility. These overlap ($r = 0.815$ at occupation level) but the 27\% of CZ-level CSH variance orthogonal to RSH does not predict polarization outcomes. This confirms the framework's scope: it measures structural feasibility, not automation displacement.


\subsubsection{Path I: Gravity Model Analysis --- COMPLETED}


\textbf{Status: COMPLETED.} Gravity model estimated on 199,362 occupation pairs. Embedding $\times$ Centroid achieves partial $R^2 = 2.55\%$; Embedding $\times$ Wasserstein achieves 3.46\%. Results consistent with \citet{CortesGallipoli2018}.

\textbf{Implication:} Task distance explains modest share of aggregate bilateral flows, consistent with structural estimates that task-specific costs are $\sim$15\% of total switching costs. Ranking divergence between conditional logit (embeddings excel) and gravity (simpler measures competitive) reflects extensive vs.\ intensive margin dynamics.


\subsection{Model Evaluation Criteria}
\label{subsec:evaluation}


The framework is evaluated against simple baselines using explicit criteria. These establish the falsifiability commitment of the research program.


\medskip
\noindent\textbf{PRIMARY CRITERIA (all must be addressed):}

\begin{enumerate}
    \item \textbf{Incremental explanatory power.} Embedding-based distance should improve fit relative to O*NET baselines by an economically meaningful margin.
    
    \item \textbf{Parameter stability.} Results should be qualitatively stable across embedding models (MPNet, MiniLM) and aggregation methods (Wasserstein, centroid).
    
    \item \textbf{Non-redundancy.} We target low-to-moderate correlation between semantic and institutional distance. With embedding Wasserstein, $\rho(d_{\mathrm{sem}}, d_{\mathrm{inst}}) = 0.25$---moderate correlation indicating the measures capture related but distinct structure.
    
    \item \textbf{Non-tautological validation.} Geometry diagnostics must use targets not used in construction.
\end{enumerate}


\medskip
\noindent\textbf{REVISION PROTOCOL:}

If primary criteria fail:
\begin{itemize}
    \item Do not add auxiliary structure to rescue fit.
    \item Report null result: ``embedding-based measure does not outperform O*NET baselines under tested specifications.''
\end{itemize}


\section{Scope Conditions and Module Limitations}
\label{sec:limitations}


This section documents scope conditions that define where the framework applies. These are features of the architecture, not failures.


\subsection{Selection Bias: Friction, Not Gates}
\label{subsec:limitation-friction}


Our institutional distance coefficient ($\gamma_{\mathrm{inst}} = 0.139$, $t = 33.7$) is estimated from completed transitions. Workers who successfully became nurses or teachers appear in our data; workers blocked by credential requirements do not. This creates selection bias \citep{Heckman1979}: we estimate the friction experienced by successful transitioners, not the blocking probability for all who attempt.

\citet{Jackson2023} finds licensing alone reduces entry probability by 24\% with no effect on exit probability---a gate structure our friction-based model cannot capture. \citet{KleinerXu2020} find licensed workers are 24\% less likely to switch occupations overall. These magnitudes suggest substantial selection: nearly one-quarter of potential transitions may be institutionally blocked and thus absent from our data. The ratio $\gamma_{\mathrm{inst}}/\gamma_{\mathrm{sem}} = 0.019$ likely understates true credential effects.

Our estimates should be interpreted as characterizing observed mobility patterns among successful transitioners, not latent worker preferences for transitions that may be prevented. Future work should incorporate exogenous credential classification (nursing, teaching, licensed trades) to model gates directly.


\subsection{Demand Side Absent}
\label{subsec:limitation-demand}


We model worker supply-side feasibility only. Job openings, employer selection, and capacity constraints are not incorporated. Predicted inflow rates exceeding 100\% for some occupations (e.g., Sales Engineers at 341\%) reflect this limitation---skill compatibility does not imply labor demand.

The supply-demand decomposition (\cref{subsec:pathway-validation}) quantifies this limitation: demand-only $\rho = 0.80$ for aggregate inflows demonstrates that supply-side feasibility alone cannot predict realized flows. Geometry provides modest information about pathway direction (per-origin $\rho \approx 0.12$) but negligible information about aggregate destinations (geometry-only $\rho \approx 0.00$). This is both a limitation and a finding---the decomposition validates the scope claim while quantifying what demand integration would add.

Feasibility is the supply-side input to equilibrium analysis. Combining feasibility with demand-side data (Lightcast vacancy data, JOLTS by occupation) would enable equilibrium predictions; we do not attempt this.


\subsection{Switching Cost Identification}
\label{subsec:limitation-costs}


Endogenous identification failed with our data. Occupation-mean wages yield $\beta_{\mathrm{wage}} < 0$, indicating that aggregate wages do not capture entry wages switchers receive. External calibration to \citet{DixCarneiro2014} is necessary, requiring maintained assumptions about cross-country and sectoral-to-occupational applicability. US estimates span 0.5$\times$ to 6.5$\times$ annual wages; our sensitivity analysis (\cref{tab:calibration-sensitivity}) shows ordinal rankings are invariant across this range while magnitudes scale proportionally.


\subsection{Marginal Distributional Contribution}
\label{subsec:limitation-distributional}


The distributional treatment of occupations as probability measures over tasks---formalized via Wasserstein optimal transport---provides theoretical grounding but no empirical advantage over simpler alternatives. After correcting for diagonal bias in the distance matrix (due to SOC$\to$Census aggregation), centroid averaging of task embeddings marginally outperforms Wasserstein while producing nearly identical distance rankings (Spearman $\rho = 0.95$) at lower computational cost. The contribution is the embedding ground metric capturing semantic task substitutability, not the optimal transport machinery per se. Future work may identify settings where distributional treatment provides greater advantage---perhaps occupations with multimodal task distributions or applications requiring the transport plan (which tasks map to which).


\subsection{Scope Delineation}
\label{subsec:limitation-scope}


This framework measures where workers \emph{can} go (structural feasibility from skill architecture), not where they \emph{do} go (realized reallocation). The former is policy-relevant for workforce development and retraining design. The latter requires general equilibrium modeling beyond current scope.

Empirically, geometry captures supply-side constraints (per-origin $\rho \approx 0.12$) while demand dominates aggregate flows ($\rho \approx 0.80$). This quantifies the scope claim: the framework measures structural feasibility---the supply-side input to equilibrium analysis. Gravity models and structural estimation can incorporate feasibility as an input to predict realized flows. The modular architecture enables this: validated $\mathcal{T}$ and $\mathcal{I}$ modules provide inputs for future $\mathcal{M}$ development.


\begin{remark}[Task distribution evolution]
\label{rem:task-evolution}
The framework's first output, $\Delta\boldsymbol{\rho}$, represents how each occupation's task distribution changes over time. This paper does not model $\Delta\boldsymbol{\rho}$ directly. We measure distances between task distributions at a point in time; we do not estimate how technology shifts which tasks an occupation uses. Modeling $\Delta\boldsymbol{\rho}$ requires either: (a) time-series data on occupation task content (e.g., from job postings), or (b) structural assumptions about how technology adoption changes task demand within occupations. This is future work. Current validation focuses on the $\mathcal{T}$ and $\mathcal{I}$ modules: measuring task-based distances and institutional barriers that determine feasibility of transitions, taking task distributions as given.
\end{remark}


\section{Applications and Research Roadmap}
\label{sec:applications}


\subsection{Theory of Change: Feasibility Under Disruption}
\label{subsec:theory-of-change}


The demand decomposition clarifies the framework's operational role. Under stable conditions, flows are demand-dominated ($\rho = 0.80$); geometry provides modest directional signal ($\rho \approx 0.12$). Under disruption---when historical demand patterns break---geometry becomes the stable input. The framework measures where workers \emph{can} go (supply-side feasibility), not where they \emph{do} go (demand-side realization). This is the operational use case: counterfactual analysis when past flows are uninformative.


\subsection{Policy Applications for Feasibility Measurement}
\label{subsec:policy-applications}


The framework supports several policy applications:

\textbf{Skill pathway analysis.} Identify ``adjacent occupations'' requiring minimal retraining. For workers in AI-exposed occupations, the geometry identifies which destinations are skill-compatible versus which require fundamental retraining.

\textbf{Workforce development targeting.} WIOA programs can prioritize structurally feasible destinations. Rather than training displaced workers for arbitrary alternatives, programs can focus on occupations within reasonable skill distance.

\textbf{Retraining program design.} Calibrated costs (8.74 wage-years per unit centroid) inform investment required. A cashier $\to$ retail sales transition costs $\sim$1.4 wage-years; a cashier $\to$ nurse transition costs $\sim$2.5 wage-years (Wasserstein-scale examples; centroid distances yield similar ordinal rankings). These planning numbers guide resource allocation.

\textbf{Green transition planning.} Map fossil fuel $\to$ renewable occupation pathways. Which occupations can absorb displaced workers from coal, oil, and gas sectors? The geometry characterizes skill-compatible destinations.

The framework does \emph{not} support: point predictions of labor reallocation, forecasts of which jobs will be created/destroyed, or capacity for absorption without demand data. The supply-demand decomposition ($\rho \approx 0.80$ demand, $\rho \approx 0.12$ geometry) quantifies why: realized flows depend heavily on labor market conditions, not just skill feasibility.


\subsection{Research Roadmap}
\label{subsec:roadmap}


The framework architecture enables incremental progress. Each phase extends one module while holding others fixed.


\begin{table}[h!]
\centering
\caption{Research Program Phases}
\label{tab:research-phases}
\begin{tabular}{llll}
\toprule
\textbf{Phase} & \textbf{Objective} & \textbf{Requirements} & \textbf{Timeline} \\
\midrule
0.8 & Demand-side integration & Lightcast/JOLTS vacancy data & 3--6 months \\
0.9 & Institutional barrier enhancement & CPS licensing supplement & 3--6 months \\
1.0 & Modality-specific shock profiles & Taxonomy design, benchmarks & 6--9 months \\
1.1 & Full GE with propagation & Individual wage data (LEHD) & 12+ months \\
\bottomrule
\end{tabular}
\end{table}


The supply-demand decomposition provides empirical grounding for Phase 0.8: demand-only $\rho = 0.80$ demonstrates that vacancy integration would substantially improve aggregate flow prediction. The modular structure is the value of infrastructure: validated components can be reused as new capabilities are added.


\section{Conclusion}
\label{sec:conclusion}


This paper proposes a modular task-space framework for analyzing labor market adjustment to technological change, and validates initial module specifications.


\textbf{Module status:}
\begin{enumerate}
    \item \textbf{$\mathcal{T}$ module (task representation):} Validated. Embedding-informed distance substantially outperforms non-embedding structured-vector baselines (74.9\% in pseudo-$R^2$). The mechanism is semantic task substitutability: sentence embeddings capture that similar tasks are interchangeable even when distinctly labeled. Wasserstein optimal transport provides theoretical grounding (minimum transformation cost) but does not outperform simpler centroid methods after diagonal bias correction; centroid-based distance is marginally preferred ($\rho = 0.95$ correlation with Wasserstein).
    
    \item \textbf{$\mathcal{I}$ module (institutional structure):} Validated. Institutional distance is associated with mobility conditional on semantic distance ($t = 33.7$). Separability holds. Directional asymmetry is sample-dependent---an informative constraint.
    
    \item \textbf{$\mathcal{S}$ module (shock profiles):} Validated. External AIOE scores are orthogonal to geometry ($r = -0.02$), confirming distinct constructs that combine coherently. Geometry substantially outperforms historical baselines on out-of-period data ($\Delta\text{LL} = +23{,}879$).
    
    \item \textbf{$\mathcal{M}$ module (mechanisms):} Preliminary. External calibration yields plausible switching costs (8.74 wage-years per unit centroid). Endogenous estimation requires individual wage data.
    
    \item \textbf{Scope:} Empirically quantified. Geometry captures supply-side pathway direction (per-origin $\rho \approx 0.12$); demand dominates aggregate flows ($\rho \approx 0.80$). Gravity model confirms task distance explains $\sim$2.5\% of aggregate bilateral flow variation, consistent with \citet{CortesGallipoli2018}. This decomposition is consistent with the framework architecture's separation of feasibility from realization.
\end{enumerate}

The framework demonstrates capability for substantive analysis: a pre/post COVID comparison finds aggregate structural stability ($\Delta\alpha < 1\%$, $p = 0.72$) with heterogeneous effects by teleworkability ($\delta_4 = -0.086$, $p = 0.01$), contributing to the debate on pandemic labor market impacts.


\textbf{Informative constraints from this phase:}
\begin{itemize}
    \item Embedding ground metric matters: semantic task substitutability drives 74.9\% improvement over O*NET baselines
    \item Distributional treatment (Wasserstein) provides theoretical grounding but no empirical advantage; centroid marginally outperforms ($\rho = 0.95$)
    \item Directional barrier patterns are sample-dependent
    \item Semantic geometry captures mobility friction, not automation susceptibility ($r = -0.06$ with RTI)
    \item Aggregate wage data insufficient for structural $\mathcal{M}$ estimation
    \item Institutional barriers estimated as friction understate gate effects
    \item Aggregate flows are demand-dominated ($\rho \approx 0.80$); geometry captures supply-side constraints ($\rho \approx 0.12$)
    \item Task distance explains modest share of bilateral flows ($\sim$2.5\%), consistent with Cortes-Gallipoli (2018)
    \item Retrospective battery confirms construct distinction: embedding geometry captures mobility friction, not automation susceptibility (Test B: 1+, 0$-$, 4(0))
\end{itemize}


\textbf{The research program continues.} The task-space framework is an organizing architecture, not a deliverable. This paper establishes that the architecture is viable: validated components ($\mathcal{T}$, $\mathcal{I}$) can be combined with external shock profiles ($\mathcal{S}$) to characterize structural feasibility. The supply-demand decomposition provides the key insight from this phase: the architecture's separation of supply-side feasibility from demand-side mechanisms is consistent with observed patterns. Geometry defines the feasible choice set; demand determines selection within that set. Future work extends $\mathcal{M}$ (equilibrium mechanisms, demand integration) and refines $\mathcal{S}$ (modality-specific profiles). The modular structure enables cumulative progress toward complete labor market modeling under technological change.


\appendix


\section{Retrospective Battery: Implementation and Results}
\label{app:retrospective}


This appendix documents the retrospective diagnostic battery testing embedding geometry against historical employment patterns. Test B (Autor-Dorn polarization) was executed. Test A (industry task drift) is blocked pending data acquisition. Test C (robot exposure) was deferred due to methodology deviation in initial implementation.


\subsection{Test A: Task Composition Shifts (ALM-Informed) --- BLOCKED}
\label{app:test-a}


\textbf{Period:} 1980--2000 (or 1980--1998 to match \citet{AutorLevyMurnane2003} exactly).

\textbf{Core question:} When computerization enters industries, do task composition shifts follow continuous distance from automated tasks, or only binary routine/non-routine classification?

\textbf{Design:}
\begin{itemize}
    \item Unit of analysis: Industry $\times$ decade
    \item Treatment: Computer capital intensity or computer use (CPS October supplements)
    \item Outcome: Change in employment-weighted task intensity by industry
\end{itemize}

\textbf{Key data requirements:}
\begin{itemize}
    \item ALM task measures (DOT 1977 based) or reconstruction via O*NET mapping
    \item Industry-level computer capital/use: BEA capital data, CPS computer use supplements
    \item Industry $\times$ occupation employment matrices: Census 1980, 1990, 2000
    \item Crosswalk: DOT 1977 $\to$ Census occupation $\to$ occ1990dd $\to$ O*NET-SOC $\to$ DWA embeddings
\end{itemize}

\textbf{Status: BLOCKED.} No industry-level data available in Dorn replication archive. Implementation requires separate ALM (2003) replication files. Deferred to future work pending data acquisition.


\subsection{Test B: Employment Reallocation (Autor-Dorn Setting)}
\label{app:test-b}


\textbf{Period:} 1980--2005 (matches \citet{AutorDorn2013} exactly).

\textbf{Core question:} Do commuting zones with higher continuous routine exposure experience different employment polarization patterns than those with equivalent discrete routine share?

\textbf{Design:}
\begin{itemize}
    \item Unit of analysis: Commuting zone ($n = 722$)
    \item Treatment (1980 baseline): Discrete RSH vs.\ continuous CSH
    \item Outcomes: $\Delta$ employment share in service occupations, $\Delta$ employment share in routine occupations, wage changes at percentiles
\end{itemize}

\textbf{Decomposition test:}
\begin{equation}
\Delta Y_{jst} = \delta_t + \beta_1 \times \text{RSH}_{j,t_0} + \beta_3 \times (\text{CSH} - \text{RSH})_{j,t_0} + X'_{j,t_0} \gamma + \phi_s + \varepsilon_{jst}
\end{equation}

If $\beta_3 > 0$ and significant, continuous captures information discrete misses.

\textbf{Replication data sources:}
\begin{itemize}
    \item Autor-Dorn replication archive: \texttt{ddorn.net/data.htm}, File [P2]
    \item Crosswalk files [A1]--[A8]: Census occupation $\to$ occ1990dd (330 balanced codes)
    \item Geographic crosswalks [E1]--[E7]: PUMA/County Group $\to$ CZ with \texttt{afactor} weights
\end{itemize}


\subsubsection{Test B Results}
\label{app:test-b-results}


\textbf{Sample:} 722 commuting zones, IPUMS Census microdata 1980--2000.

\textbf{CSH construction:} Continuous Semantic Height (CSH) is the projection of occupation embedding centroids onto the learned RTI direction in embedding space. Ridge regression ($\alpha = 596.4$, 5-fold CV) prevents overfitting when learning the RTI direction from 16 O*NET ability elements. At the occupation level, $r(\text{CSH}, \text{RTI}) = 0.815$, confirming the embedding geometry captures economically meaningful task structure.

\textbf{CZ aggregation:} Employment-weighted average of occupation-level CSH values. At the CZ level, $r(\text{CSH}_{\text{cz}}, \text{RSH}_{\text{cz}}) = 0.478$. The decorrelation from 0.815 to 0.478 reflects aggregation across different occupation compositions---CZs with the same RSH can have different CSH depending on which specific routine occupations are present.

\textbf{Specification:} State fixed effects, year fixed effects, population weights. CSH residualized against RSH to isolate orthogonal information: $\text{CSH}_{\text{resid}} = \text{CSH}_{\text{cz}} - (\alpha + \beta \cdot \text{RSH}_{\text{cz}})$. Residual variance retained: 73.2\%.


\begin{table}[h!]
\centering
\caption{Test B Results: CSH Incremental Explanatory Power}
\label{tab:test-b-results}
\begin{tabular}{lccccl}
\toprule
\textbf{Outcome} & $\beta(\text{CSH}_{\text{resid}})$ & \textbf{SE} & \textbf{$p$-value} & $\Delta R^2$ & \textbf{Verdict} \\
\midrule
$\Delta$ operator share & $-3.337$ & 0.342 & $<0.001$ & 5.5\% & + \\
$\Delta$ routine share & $-1.034$ & 0.315 & 0.001 & 0.77\% & 0 \\
$\Delta$ clerical/retail & $+0.437$ & 0.331 & 0.187 & 0.09\% & 0 \\
$\Delta$ service share & $+0.214$ & 0.364 & 0.557 & 0.02\% & 0 \\
$\Delta$ mgmt/prof/tech & $+0.545$ & 0.311 & 0.079 & 0.13\% & 0 \\
\bottomrule
\end{tabular}
\end{table}

\textbf{Verdict thresholds:} ``+'' requires $p < 0.05$ AND $\Delta R^2 \geq 1\%$.

\textbf{Summary:} 1+, 0$-$, 4(0). Embedding geometry adds modest incremental explanatory power for operator employment decline ($\Delta R^2 = 5.5\%$), but the orthogonal component of CSH does not systematically predict other polarization outcomes.

\textbf{Interpretation:} The weak incremental explanatory power is expected given construct distinction. CSH captures task similarity and mobility friction---where workers \emph{can} go given skill requirements. RTI captures automation susceptibility---which jobs technology will displace. These constructs overlap substantially ($r = 0.815$ at occupation level) because routine tasks are both automatable and have similar skill requirements. However, the 27\% of CZ-level CSH variance orthogonal to RSH does not predict employment polarization. This clarifies scope: the framework measures structural feasibility, not automation displacement.

The operator exception ($\Delta R^2 = 5.5\%$) may reflect meaningful gradations within routine manual work---continuous distance captures variation that binary classification misses. Alternatively, with five outcomes tested at $\alpha = 0.05$, one significant result is consistent with chance. We treat this as suggestive, not confirmatory.


\subsection{Test C: Robot Displacement (Acemoglu-Restrepo Setting) --- DEFERRED}
\label{app:test-c}


\textbf{Period:} 1990--2007 (or 1993--2014 depending on IFR data availability).

\textbf{Core question:} Does continuous task overlap with robot-exposed activities perform better than binary robot exposure classification for probabilistic scoring of employment/wage changes?

\textbf{Design:}
\begin{itemize}
    \item Unit of analysis: Commuting zone or occupation
    \item Treatment: Industry-level robot exposure (A-R) vs.\ task-level robot exposure (ours)
    \item Outcomes: Employment-to-population ratio, log wages
\end{itemize}

\textbf{Data sources:}
\begin{itemize}
    \item IFR robot data by industry $\times$ country (1993--2014)
    \item A-R replication files
    \item European robot adoption for IV
\end{itemize}

\textbf{Status: DEFERRED.} Initial implementation used keyword classification of robot-relevant DWAs rather than embedding distance to robot-task centroid as specified. This methodology deviation means the implemented test validates a classification scheme, not embedding geometry. Correct implementation (embedding-based robot exposure) deferred to future work.


\subsection{Interpretation Matrix}


\begin{table}[h!]
\centering
\caption{Retrospective Battery Status Summary}
\label{tab:battery-summary}
\begin{tabular}{llll}
\toprule
\textbf{Test} & \textbf{Status} & \textbf{Key Result} & \textbf{Interpretation} \\
\midrule
B (Polarization) & Executed & 1+, 0$-$, 4(0) & Construct distinction confirmed \\
C (Robots) & Deferred & Methodology deviation & Requires re-implementation \\
A (Task Drift) & Blocked & No industry data & Requires ALM files \\
\bottomrule
\end{tabular}
\end{table}


\textbf{Framework implication:} The retrospective battery confirms the framework's scope. Embedding geometry captures mobility friction (validated via CPS transitions in \cref{sec:validation}) but does not substantially improve automation prediction beyond discrete RTI measures. This is consistent with the paper's claim that geometry measures where workers \emph{can} go, not which jobs will be displaced.


\begin{table}[h!]
\centering
\begin{tabular}{|l|c|c|c|l|}
\hline
\textbf{Pattern} & \textbf{Test A} & \textbf{Test B} & \textbf{Test C} & \textbf{Implication} \\
\hline
Continuous dominates & $+$ & $+$ & $+$ & H1: Manifold approach preferred throughout \\
\hline
Discrete dominates & $-$ & $-$ & $-$ & H2: Discrete classification sufficient \\
\hline
Adjustment only & $?$ & $+$ & $-$ & H3: Continuous for reallocation, discrete for displacement \\
\hline
Noise & $0$ & $0$ & $0$ & Semantic similarity $\neq$ economic substitutability \\
\hline
\end{tabular}
\caption{Interpretation matrix for retrospective diagnostic battery. ``$+$'' = continuous outperforms discrete; ``$-$'' = discrete outperforms; ``$0$'' = no significant difference.}
\label{tab:interpretation-matrix}
\end{table}

\textbf{Observed pattern:} Test B yields 1+, 4(0)---weak support for continuous measures in one outcome (operator decline), no difference in four others. This is most consistent with the ``Noise'' row interpretation for automation prediction, while mobility prediction (CPS validation) shows strong continuous advantage. The framework is validated for its intended purpose (mobility/feasibility) but does not extend to automation prediction.


\bibliographystyle{apalike}
\bibliography{references}


\end{document}
